% T7 of the Transport Papers series (consolidated).
% Correction (2026-09-10, mechanical pass, per Leon): rule 23 misidentified T5 as the printed-version exception. T1 is the exception at Version 3.0; every other paper, T5 included, reads Version 1.0. \date corrected accordingly (was 1.1).
% v1.1 (2026-09-10, derivation session, Fable 5.1 subagent, per rule 20,
%   series-review repair pass against series_review_2026-09-09.md).
%   The review found nothing in T7's body or script. Two data items only:
%   \bibitem{ButlerT3} given T3's full subtitle (Iterated Transport, the
%   Frozen Clock, and Second Quantisation), matching the other papers;
%   \date -> Version 1.1. No body text, label, number or claim changed.
%   Script: verify_T07.py 25/25, re-run.
%   Compile: pdflatex x2, 21 pages, 0 errors, 0 undefined refs, 0 multiply
%   defined, 0 overfull, no float overflow.
% 2026-09-09 addendum (interactive session with Leon, Opus 5). Two changes
%   per Leon: (1) the printed version is now 1.0 for every paper T0 to T17,
%   T5 alone at 3.0 (progress-doc rule 23); the development numbering is
%   retired from the title page and kept in this header. (2) Every reference
%   to the programme monograph is removed (rule 24), re-pointed to the series
%   where a sibling paper carries the result and otherwise restated to stand
%   alone. No mathematics, no labels and no numbering changed, so verify_T07.py
%   stands unchanged at 25/25 and was not re-run.
%   Here: version only. NOTE: this file was restored from the project copy in
%   the same session after an editing error truncated it; the restore was
%   verified by structure and a clean compile at 21pp.
% v1.8 (2026-09-09, derivation session, Fable 5.1 subagent, per rule 20,
%   review-repair pass against T07_review_v6.md, same session). The review of
%   the v1.6 prose pass and the v1.7 successor problem 6 found 0 BLOCKING and
%   0 SERIOUS findings; the cheap NON-BLOCKING ones applied: N1 Rem 6.4
%   "the series quantifies its own claims" -> "states its own claims at their
%   level" (the glossary's wrong alternative for "prices"); N2 Assembly 5.3
%   Joint J3 "reconstruction composed with the transform" -> "reconstruction,
%   each helicity sector identified through the transform of Joint J1" (no
%   composition-of-maps claim at a modelling joint); N3 successor problem 6
%   "T8 currently takes" -> "T8 takes"; N4 intro fourth question "this
%   ell_P^2/A correction" -> "the resulting ell_P^2/A correction"; N5 script
%   label for the homogeneity table d -> h (the paper's symbol since v1.3);
%   N6 script docstring's pre-v1.4 subtitle replaced by the main title alone,
%   as the ButlerVerifyT7 bibitem quotes it. NOT changed, for Leon: N8 the
%   title line's "Ledger" and "Finite-Ledger Seam" (his v1.4 ruling); N7 the
%   pre-existing literal homogeneity-table check; N9 the abstract's
%   "computational" trim; N10 recorded as a scope narrowing, not a defect.
%   Check count unchanged, 25/25 (labels only). Compiles clean x2: 0 errors,
%   0 undefined, 0 multiply-defined, 0 overfull, no float overflow, 21pp.
%   Numbering frozen: aux diff against v1.7 empty. KNOCK-ON noticed, not
%   applied: T6 (line 923), T8 (985), T10 (1093), T12 (1666) bibitem ButlerT7
%   print the pre-v1.4 subtitle "State, Continuum, and the Four-Dimensional
%   Lift"; T9 and T17 print the current one.
% v1.7 (2026-09-09, derivation session, Fable 5.1 subagent, per rule 20, gap
%   merge pass: successor problem T7 never recorded). NEW successor problem 6,
%   appended at the end of the Section 8 enumerate: the helicity-counting
%   question (does one boundary mode of the near-horizon theory carry one
%   helicity channel of the lift, or both helicities of a massless species?),
%   its downstream stakes in T8 (capacity per boundary mode ln 2 against
%   2 ln 2, kappa_c = 1/(4 ln 2) against 1/(8 ln 2), cited [Prop. 4.12]{ButlerT8},
%   number and kind checked against the suite aux, prop:residue 4.12), its
%   location in this paper (joint J3 of Assembly asm:helicity, the modelling
%   joint; J1 and J2 place the helicities but do not count channels per mode),
%   and the note that T8 takes the one-helicity reading at the tier of that
%   joint. Existing items 1-5 untouched. Count "Five problems" -> "Six
%   problems" (the only place the count is printed; abstract, intro, levels
%   list and Position do not quote it). \date bumped 1.4 -> 1.7 (it had been
%   left at 1.4 by earlier passes). ButlerT8 bibitem already present, no
%   change. No script change: no check quotes the count. Compiles clean x2:
%   0 errors, 0 undefined refs, 0 multiply-defined, 0 overfull, no float
%   overflow, 21pp. verify_T07.py 25/25 (unchanged). Numbering frozen: aux
%   diff empty (enumerate items are unlabelled). No knock-ons noticed.
% v1.6 (2026-09-08, prose pass, Fable 5.1 subagent, per project rule 22: mannered
%   prose replaced by literal statement, Leon's ruling to replace the recurring
%   vocabulary too). Body only, from \begin{document}; preamble, title line and
%   bibliography untouched; no mathematics, labels, numbering or claims changed.
%   Compiles clean x2: 0 errors, 0 undefined refs, 0 multiply-defined, 0 overfull;
%   21pp (was 20: the last three bibliography lines now spill to a 21st page).
%   verify_T07.py 25/25 (unchanged; script untouched). Numbering unchanged
%   (aux diff empty). Abstract 250 words. Glossary applied: "ledger" (black-hole
%   sense) -> "boundary mode space"/"boundary modes" (about 30; "ledger algebra"
%   -> "boundary-mode algebra", "continuum/countable ledger" -> "continuum/
%   countable mode space", "finite-ledger" -> "finite-mode"); "ordering rent" ->
%   "ordering shift" (12); "seam" -> "the ell_P^2/A correction" (10, incl. Sec 6
%   title, Rem 6.3 title, successor problem 4; the label sec:seam/rem:seam stays);
%   "priced/prices" -> "stated/states/quantifies" (5); "honest/honestly/honesty"
%   dropped or -> "correct/declared/at what level" (6); "touchpoint" -> "point of
%   contact" (3); "owned by" -> "assigned to"; "forward ledger" -> "list of open
%   directions"; "series ledger of T17" -> "claim table of T17"; "bookkeeping/
%   machinery" -> "results/result" (5); "business" -> "task"; "lives/live in" ->
%   "is defined in/are used inside" (4); "face" -> "instance"; "engine" ->
%   "arithmetic". Non-glossary: "frontier" -> "open problem/problem" (4, label
%   rem:frontier stays); "algebraic home" -> "the algebra carrying" (2);
%   "custody" -> "shared/question shared"; "houses/housed/housing" ->
%   "contains/has a canonical place/location"; "spine" -> "core"; "mapped this
%   ground" -> "treated this question"; "switch off" -> "vanish"; "the ultraviolet
%   story" -> "behaviour"; "easier to hold than to test" -> "at present untestable";
%   "the debts are declared" -> "the prior work is acknowledged".
% 2026-08-24: renumbered (was T9) and retitled per Leon; content unchanged.
% v1.0 (sprint 2, 2026-08-21): review draft, M1-M6 complete (state input,
% mode continuum, four-dimensional lift, finite-ledger seam, scope,
% verification 22/22; rule-13 figure verdict: zero figures).

% v1.1 (2026-09-01, review pass on Leon's T07_review_v2.md; derivation session):
%   ALL 10 BLOCKING, 18 SERIOUS and 22 MINOR findings addressed, one deliberately
%   left (below). Compiles clean x2: 0 errors, 0 undefined refs/citations,
%   0 overfull boxes, 19pp. verify_T07.py: 26/26 pass (was 22; the hollow
%   thm:istate(1) check replaced by two genuine ones, re-hardened by mutation:
%   the v2 review's non-thermal-weight mutation now fails both). Suite
%   crossrefs re-run: 377/377, 0 broken. Numbering: prop:mode 2.1,
%   fact:quasifree 3.3, rem:whichdynamics 4.5, prop:edge 6.1 all PRESERVED, so
%   no inbound locator in T8/T10 moved; T12's [Prop. 5.2] now lands on the new
%   homogeneity proposition, deliberately.
%   The mathematics: B1 sigma = e0e1 in the statement (sign); B2 the kappa
%   convention settled (omega_j = Schwarzschild frequency, h = diag(omega_j),
%   beta = 2pi/kappa; rapidity bridge sentence in Def 2.7); B3 the garbled
%   prop:mode(2) proof rewritten; NEW Prop 5.2 (prop:homogeneity) DERIVES the
%   Penrose homogeneity -2s-2 from T9's ordering rent (q(id) = -hbar(Upsilon+2),
%   script-checked on monomials), discharging T9's successor problem 4, per
%   Leon's ruling; NEW Thm 3.6 clause (2): the KMS identity forces the pairing
%   to vanish (gauge invariance derived, S2), so the theorem now covers the
%   whole quasi-free class; pure-import clauses moved out to NEW Cor 3.8;
%   thm:continuum restated for self-adjoint h (S7), fixing the Sec 6 hypothesis
%   clash via the bilateral-ladder bridge; prop:edge gains the explicit
%   cross-commutator (S9), the 2m <= L clause (M3) and an index argument in
%   place of "telescopes" (S14); prop:planck tail-bound chain repaired (S17),
%   window sum defined (M4), m-times-BE normalisation stated (M5), (Q)-ladder
%   flagged as modelling (S6).
%   Provenance: Fact 2.2 rewritten against T5 v1.19 (false converse REMOVED,
%   premises (P)/(R) named, beta-free-in-sign stated; B4); Fact 2.6 shrunk to
%   the rent T9 proves (B5); Fact 2.4 split unconditional/(Q)-conditional with
%   Prem 8.1 + Thm 8.2 locators (S6); Fact 2.5 scoped to Lambda = 0 (M13);
%   WW triage de-attributed from T5, owned at Tier C (B7); monograph bibitem
%   REMOVED and all nine prose mentions owned or dropped (B9, rule 17b); NEW
%   Facts 3.4/3.5 for the BR 5.3.25 and Araki-reduction imports (S5).
%   Structure/house rules: helicity lift and greybody now in a NEW Assembly
%   environment, no proof env; interface premise (I) lifted into the greybody
%   statement (S3, per Leon's route-1 ruling); NEW Position-in-the-series
%   section (B8, rules 15/18) with bibitems ButlerT0/T8/T10/T12/T17 added;
%   six-step chain and premised thermality per T5 v1.5 (B6), novelty re-priced
%   to the operator-level all-modes statement (abstract + intro rewritten,
%   Tier A/B/C vocabulary, four paragraphs, contribution enumerated);
%   rem:frontier records T5's (G)-rejection and rate rigidity (S16); units
%   sentence (S15); notation: L for the mode count, u/w scalars, cal-N flux,
%   d for homogeneity degree (S12); unused packages dropped (M21).
%   DELIBERATELY LEFT, PER LEON'S RULING (2026-09-01): the two "XV" references
%   (rem:liftscope and successor problem 2). FLAG FOR DECISIONS NEEDED: the
%   review's T13/T14 proposal conflicts with the roadmap's "T10 (fuzzy sphere)"
%   dependency tag, the progress doc's T8-M3 "shared custody with T10 named",
%   and the 08-24 mapping (original T15 = T10); resolve the custodian, then
%   replace both XVs and re-run crossrefs.
%   KNOCK-ONS APPLIED THE SAME SESSION: T12 fact:helicity reworded (rent in
%   T9, derivation in T7; locator unchanged and now true); T17 T5 row
%   six-step, T7 row 19pp/26 with re-priced payload, series total 987 -> 991
%   in prose + verify_T17.py (audit 37/37). Both recompiled clean.
% v1.2 (2026-09-01, same day, from Leon's v3 review of v1.1): both blocking
%   items were v1.1's own defects, owned. (1) The successor-section splice had
%   left \section{Machine verification} DUPLICATED on one line (an empty Sec 9
%   before the real one); compile-clean checks cannot see it; removed. (2) The
%   v1.1 pairing check was a no-op (asserted 1 + x_j != 0), the very hollow
%   pattern it was written to remove. REPLACED by a real solve on a two-mode
%   window: the full KMS constraint system for an arbitrary density matrix
%   (16 unknowns) against independently supplied dynamics; kernel confirmed
%   one-dimensional, spanned by the Gibbs matrix, unique solution confirmed
%   to have zero pairing and the FD symbol via Jordan-Wigner matrices.
%   Mutation-hardened: a non-thermal dynamics weight and a constraint sign
%   flip both FAIL it (a JW-sign drop passes, correctly: the asserted
%   quantities are sign-blind on the diagonal state). Still 26 checks.
%   Also from the v3 review: prior-art sentence and Contribution (1) re-priced
%   to match the abstract (T5's single-mode derivation credited); the -2s-2
%   promotion stated honestly everywhere (the computation appears inside T9's
%   proof of Thm 5.7 and is checked by T9's script; this paper promotes it to
%   a stated result, discharging T9's successor 4), wording changed in the
%   abstract, intro, contribution, prior art and Prop 5.2's proof; tier
%   scaffolding reconciled both ways (abstract now files the promotion under
%   Tier B's proved joint and conditions the Planck-sector Tier A claim on the
%   (Q)-convention ladder; the ledger's Prop 6.2 row names the ladder).
%   XV RESOLVED per Leon's ruling (2026-09-01): XV = T10, on the roadmap's
%   "T10 (fuzzy sphere)" dependency tag, the roadmap's own T7 entry, and the
%   progress doc's T8-M3 custody note; both occurrences now read T10. The
%   v1.1 header's renumbering rationale (original T15 = T10) was wrong per
%   Leon and is withdrawn; the ruling rests on the records above.
%   KNOCK-ONS, same session: T8 fact:inventory now carries Prop 2.1(5)'s
%   up-to-reparametrisation qualification; T9 successor problem 4 records its
%   discharged half citing [Prop. 5.2]{ButlerT7} (bibitem added), transform
%   half kept open; T17 caption and verify_T17.py docstring 534 -> 991. All
%   three recompiled clean; verify_T08/T09/T17 all pass (37/37 audit).
%   This version: compiles clean x2, 0 errors, 0 undefined, 0 overfull, 20pp
%   (originally recorded as 19pp in error; corrected 2026-09-01, v1.3 pass);
%   verify_T07.py 26/26; suite crossrefs re-run after the T9 edit.
% v1.3 (2026-09-01, same day, from Leon's v4 review of v1.2; this pass worked
%   the FULL review document, not the message summary, closing the process gap
%   the v4 verdict named). BLOCKING: (1) Prop 6.2(3) said the window sum
%   exceeds the bilateral sum; it is the smaller, as its own proof and the
%   coverage text say; statement now reads "the bilateral sum exceeds the
%   window sum by a positive amount bounded by...". (2) Fact 2.6 imported the
%   ordering shift unsigned (+hbar/2 per dimension); it is -hbar/2, totalling
%   -2 hbar, the sign that makes q(id) = -hbar(Upsilon+2) and hence -2s-2;
%   the Fact now carries the sign and T9's closed form q(M) = Z M Zbar -
%   (hbar/2) tr M, and Prop 5.2's statement sign matches. TIERS: abstract
%   rewritten in section order with the edge identity unconditional, the
%   Planck factor under (Q), full-algebra existence credited to the import,
%   and the l_P^2/A reading moved to modelling level (250 words); bridge
%   paragraph at the head of Sec 7 maps the ledger onto Tier A/B/C and names
%   the paper's own tier (B, per T17); Sec 5 opener tagged Tier B and now
%   counts the greybody premise (I) alongside the three joints. PROVENANCE:
%   discharge language halved everywhere ("the homogeneity half...; the
%   transform itself stays open"), matching T9; Fact 2.5's scattering
%   sentence (T7's own inference) deleted, its content living in the greybody
%   assembly's joint (i), and its T6 locator extended to Thms 4.5 and 4.7;
%   Fact 2.6's provenance extended to [Thm. 3.1, Rem. 3.2, S 5.2, Lem. 5.5,
%   Thm. 5.7] with the promoted-from-the-proof clause; rem:frontier's bare
%   T5 locators now a proper cite; Contributions (1)/(3) cite Thm 7.20 and
%   Thm 5.7; Prop 5.2's provenance commentary moved from proof to statement.
%   CLARITY: twistor indices renamed mu/nu (no more alpha/beta collision with
%   the modular objects); homogeneity degree d -> sans-serif h; b_m -> hat
%   a_m with the bilateral-vs-finite-window citation made honest; horizon
%   area bound to A; l_P exemption stated; cal-N glossed; truncation J -> j_0;
%   WW closing sentence rewritten (no more symmetric "as hard to confirm as
%   refute" calling itself an asymmetry); assembly labels prop: -> asm:;
%   Thms 3.6(3)/4.2(3) statements aligned to the proved operator order;
%   unbounded-h extension step fixed (bounded identity extends, unbounded
%   stays on the core); core appositive fixed (the UNION is the core);
%   Planck tail bounds now valid for every index (no case split, covers
%   K < m); cross-commutator m < m' by adjoints, said; beta in R in
%   lem:gibbs; intro's fourth question got its own paragraph; Position in
%   T1's two-paragraph shape; storey glossed; example/conjecture/question
%   envs dropped; verb -> texttt; T8 epithet unified ("the entropy paper");
%   (2)--(3) en dashes; verification coverage REWRITTEN as a per-section
%   description list in document order, Sec 4 accounted for, the duplicate
%   Fermi-Dirac detailed-balance check REMOVED from the script (25 checks
%   now, honestly counted; the Bose-Einstein form retagged to prop:planck(2)),
%   and the rotor prose corrected (R = 5/3 + (4/3) sigma, cosh(lambda/2) =
%   5/3). Ledger successors now an enumerate.
%   KNOCK-ONS, same session: T9's uncited monograph bibitem removed (rule
%   17b); T8's ButlerT10 bibitem normalised ("forthcoming" dropped); T17 T7
%   row 20pp/25 with "promoted", series totals 991 -> 990 everywhere
%   including verify_T17.py (audit 37/37). T8 and T9 recompiled clean; T17
%   recompiled, its three pre-existing "Float too large for page" warnings
%   unchanged (v5 review, sibling item 2).
%   AWAITING LEON'S RULING (not edited): the subtitle's omission of the seam
%   (v4 finding 29); T17's "preprint of record" monograph citation (v4
%   sibling item 3: possible standing exemption for the survey, unrecorded).
%   This version: compiles clean x2, 0 errors, 0 undefined, 0 overfull, 20pp;
%   verify_T07.py 25/25; suite crossrefs re-run.
% v1.5 (2026-09-01, from Leon's v5 review of v1.3/v1.4; all twenty findings
%   applied, in T7 and in one sibling). THE FIVE SERIOUS, one clause each:
%   (1) the abstract's "Tier A, proved here:" headed a clause list that then
%   twice said "imported"; ", proved here" deleted, the per-clause wording
%   already pricing each item. (2) The abstract detached the window tails from
%   premise (Q), which Prop 6.2's standing hypothesis and the ledger both
%   attach; the tails now sit with the Planck factor on the (Q) ladder, the
%   edge identity alone unconditional. NOTE: that wording came from the v4
%   review's own suggested fix, which was right about Prop 6.1 and wrong about
%   Prop 6.2(3); recorded so the error is not reintroduced. (3) The Sec 7 tier
%   bridge mapped Prop 5.2 to Tier A where the abstract meets it at Tier B and
%   fused the Tier B and modelling buckets the abstract had just separated;
%   the paragraph now DEFINES the three tiers, names Prop 5.2 explicitly, and
%   keeps assembly and modelling distinct. (4) Fact 2.2 called 2pi/kappa a
%   rapidity period; the rapidity period is 2pi and 2pi/kappa is Schwarzschild
%   time, as Def 2.7 and the Fact's own next clause both say. This survived
%   four rounds because the v4 review filed it under "Verified sound" rather
%   than as a numbered finding. (5) Prop 5.2's proof identified the commutator
%   sum (-hbar per dimension) with Fact 2.6's shift (-hbar/2 per dimension);
%   the halving by symmetrisation is now said. The v1.3 sign fix exposed it.
%   THE FIFTEEN MINOR: "whose one proved joint" -> "whose only joint proved in
%   this paper"; T17 cited in the bridge; Fact 2.6's provenance corrected to
%   [Thm. 3.1, Rem. 3.2, S 5, Lem. 5.5, Prop. 5.6, Thm. 5.7] (the palatial
%   relation is in T9's Sec 5 preamble, not S 5.2; Prop 5.6 is the Clifford
%   side of the rent); the coverage list's Bose-Einstein form corrected to
%   nbar/(1+nbar) = e^{-beta eps m} (N was the number operator, and the
%   exponent was the wrong variable), and the script's own label for that
%   check renamed N -> nbar to match; the horizon area A bound at its first
%   use in the units sentence; a sans-serif hedge for the homogeneity degree
%   against the one-particle Hamiltonian h; Sec 5's opener now says the
%   HELICITY assembly's joints and records that greybody itemises three more;
%   the coverage list's section numbers driven by \ref; the successor-2
%   epithets ("the entropy paper T8", "the action paper T10"); T1's italic
%   leads in the Position section; "storey" glossed at its first use in the
%   intro; two sites saying "inside the dictionary theorem" -> "inside the
%   proof of"; the index parenthetical restated as a convention rather than an
%   announcement; long single-sentence coverage items and ledger paragraphs
%   split. Abstract retrimmed to 247 words (was 252 after the fixes, over the
%   rule-8 ceiling): "there is exactly one that is KMS" -> "exactly one is
%   KMS", and the open-problems list de-articled.
%   NO result renumbers: the freeze of v1.4 holds, and no inbound locator in
%   T8/T9/T10/T12 moves. verify_T07.py UNTOUCHED at 25/25 (re-run in Claude's
%   workspace, not on this machine, per Leon's standing instruction).
%   KNOCK-ON, same session: T8's uncited \bibitem{Butler2026} removed (rule
%   17(b); the same defect this session's v1.3 pass removed from T9, three
%   lines from the bibitem it did edit), T8 recompiled clean. The roadmap's T7
%   entry gained its missing T6 dependency, in both the Depends-on line and
%   the dependency graph; T7's and T6's Position sections already agreed with
%   each other and the roadmap was the stale record. NOT DONE, and left for
%   Leon: the same uncited monograph bibitem is live in T1, T2, T3, T4, T5, T6,
%   T10 and openitems_draft.tex, several of them REVIEWED & DONE with frozen
%   numbering, so whether rule 17(b) is applied per-touch or in one sweep is
%   his ruling; T17's "preprint of record" (v4 sibling item 3) is the same
%   question and also stands.
%   This version: compiles clean x2, 0 errors, 0 undefined, 0 overfull, 20pp.

% v1.4 (2026-09-01, per Leon's ruling on the v4 review's finding 29): subtitle
%   now names the seam: "State, Continuum, Lift, and the Finite-Ledger Seam".
%   No other change. Compiles clean x2, 0 errors/undefined/overfull, 20pp;
%   verify_T07.py 25/25. REVIEWED & DONE per Leon, 2026-09-01: four review
%   rounds (his review docs v1-v4 in the project; T07_review_v1/v2 also in the
%   folder), all findings closed or ruled on; numbering FROZEN at this version
%   per rule 21; verify_crossrefs run at declare-done. Second open ruling from
%   the v4 review (T17's "preprint of record" monograph citation, sibling item
%   3) remains with Leon and is not a T7 matter.
% House rules: no em dashes; short abstract; standard series disclaimer;
% physics touchpoints kept (rule 10); proof standard per rule 11 (Fact
% (imported) discipline); verification imports GA.py (rule 12); diagrams
% per rule 13.
\documentclass[11pt]{article}
\usepackage[margin=1.1in]{geometry}
\usepackage{amsmath,amssymb,amsthm}
\usepackage[dvipsnames]{xcolor}
\usepackage[colorlinks=true,linkcolor=MidnightBlue,citecolor=MidnightBlue,urlcolor=MidnightBlue]{hyperref}

\newtheorem{theorem}{Theorem}[section]
\newtheorem{proposition}[theorem]{Proposition}
\newtheorem{lemma}[theorem]{Lemma}
\newtheorem{corollary}[theorem]{Corollary}
\theoremstyle{definition}
\newtheorem{definition}[theorem]{Definition}
\newtheorem{remark}[theorem]{Remark}
\newtheorem{fact}[theorem]{Fact}
\newtheorem{assembly}[theorem]{Assembly}

\newcommand{\Cl}{\mathrm{Cl}}
\newcommand{\R}{\mathbb{R}}
\newcommand{\C}{\mathbb{C}}
\newcommand{\Z}{\mathbb{Z}}
\newcommand{\Mat}{\mathrm{M}}
\newcommand{\CAR}{\mathrm{CAR}}
\newcommand{\End}{\operatorname{End}}
\newcommand{\tr}{\operatorname{tr}}
\newcommand{\id}{\mathrm{id}}
\renewcommand{\TH}{T_{\mathrm{H}}} % \renewcommand: LaTeX predefines \TH as the thorn glyph (the T5 lesson)
\newcommand{\ip}[2]{\langle #1, #2\rangle}
\newcommand{\T}{\mathbb{T}}

\title{\bfseries Hawking Radiation from the Ledger\\[4pt]
  \large\itshape State, Continuum, Lift, and the Finite-Ledger Seam}
\author{Leon Butler\\ \normalsize\href{mailto:lbutler.math@gmail.com}{\texttt{lbutler.math@gmail.com}}\\ \normalsize ORCID: \href{https://orcid.org/0009-0006-4449-5617}{0009-0006-4449-5617}}
\date{Version 1.0, September 2026 \textperiodcentered{} Transport Papers, T7}

\begin{document}
\maketitle

\begin{abstract}
\normalsize
This paper is a quantum companion to T5: it supplies the state, the continuum, and the four-dimensional lift between T5's single-mode core and the physical Hawking spectrum, each claim at its tier. Tier A: among the quasi-free states of the boundary modes' CAR algebra exactly one is KMS for the geometry's boost at the Hawking temperature, the I-state; existence is window-exact, carried to the full algebra by imported inductive-limit results; the KMS identity itself forces the pairing to vanish and the symbol to be Fermi-Dirac, all modes at once, extending T5's single-mode occupancy theorem; uniqueness beyond the class is the imported ideal-Fermi-gas theorem; and the statements carry to the continuum on a dense core. Tier B, the joints itemised: the helicity lift and greybody inheritance, whose only joint proved in this paper, the homogeneity $-2s-2$, is promoted to a stated result from inside the proof of the twistor paper's dictionary theorem; the exterior coincides with general relativity on the invertible region at $\Lambda = 0$, so the flux prediction equals Hawking's, filtering included. Tier A again: the finite mode space's deviations from the free-boson current algebra are edge-supported by an exact commutator identity; and, on the linear s-wave ladder of T5's premise (Q), an exact Planck factor with proven exponential window tails. The $\ell_P^2/A$ scale of generic corrections is a dimensional reading at modelling level; the per-observable scaling is open. Open and named: axiomatic selection of the dynamics, the lift beyond assembly, the Kerr continuum. Exact-arithmetic checks verify the identities.
\end{abstract}

\medskip
\noindent\textbf{MSC 2020.} Primary 81T05, 46L60; Secondary 83C57, 46L30, 81R10.\\
\textbf{Keywords.} Hawking radiation; KMS states; quasi-free states; CAR algebra; modular theory; thermofield double; Penrose transform; greybody factors; bosonisation.

\newpage
\textbf{Authorship and disclosure. }The founding idea of the series is the author's: that division by zero in a Clifford algebra should be read not as a prohibition but as a transport into a higher-dimensional algebra, and that the transported division need not be accompanied by an inverse operation. The direction of investigation and the critical review of the results are likewise the author's. The mathematical development and the drafting of the text were carried out in extended collaboration with Claude Fable 5 (Anthropic), an AI system.

\newpage
\section{Introduction}\label{sec:intro}

T5 of this series \cite{ButlerT5} computes a six-step chain from the Schwarzschild gauge field to the Page curve. Two of its links are deliberately minimal, and say so: single-mode thermality is proved from two named premises, (P) univalence and (R) the KMS condition at the geometric period, with the occupancy derived from the gluing identity \cite[Prems.~7.8 and 7.19, Thm.~7.20]{ButlerT5}; and the bosonic Planck sector is derived in the s-wave-reduced theory under a joint quasi-freeness premise \cite[Prem.~8.1]{ButlerT5}. Between that core and the physical Hawking spectrum stand four questions, and this paper answers them at their correct levels.

\emph{Which state?} Section \ref{sec:state} proves the I-state theorem, the I-state being the series' name for the framework's distinguished equilibrium state $\omega_I$: among the quasi-free states of the CAR algebra of the boundary modes there is exactly one that is KMS for the geometry's boost at the Hawking temperature, and the KMS identity forces its two-point data outright, the pairing to zero and the symbol to Fermi-Dirac (Theorem \ref{thm:istate}). T5 derives the single-mode occupancy the same way; the increment here is operator-level and simultaneous across all modes, with uniqueness inside the quasi-free class proved, and uniqueness beyond it imported with provenance. What is genuinely open is the \emph{selection} of the dynamics from a geometric axiom, and Section \ref{sec:state} poses that open problem precisely, in the terms T5's own selection analysis fixes.

\emph{Which algebra carries the mode continuum?} Section \ref{sec:continuum} extends the single-mode statements over the tower's CAR limit \cite{ButlerT3}: the finite-window statements are exact theorems, the inductive-limit results are imported as Facts with provenance, and the result is free-fermion thermal field theory as the algebra carrying the near-horizon boundary modes.

\emph{How does the two-dimensional core become four-dimensional radiation?} Section \ref{sec:lift} assembles the lift with its joints itemised: helicity is the charge of the twistor storey (one level of the completion tower), and the Penrose homogeneity $-2s-2$ is proved from the ordering shift of the twistor paper \cite[Thm.~5.7]{ButlerT9} in Proposition \ref{prop:homogeneity} (the computation appears inside that theorem's proof; the proposition promotes it to a stated result); greybody factors are inherited verbatim, because the series' exterior description coincides with general relativity on the invertible region at $\Lambda = 0$ \cite{ButlerT6}, so the flux prediction equals Hawking's, filtering included.

\emph{What does the finite boundary mode space change?} Section \ref{sec:seam} proves that the deviations of the finite boundary mode space from the infinite current algebra are edge-supported, with the exact edge identity and quantified tails: the corrections are area-suppressed at their proven rates, and the resulting $\ell_P^2/A$ correction is stated as a prediction. The Weinberg-Witten triage against the composite-graviton reading is recorded there as this paper's own claim.

\subsection*{Prior art}

The prior work is acknowledged once, here. Hawking's radiance is \cite{Hawking1975}. The thermofield-double reading of the eternal black hole is Israel's \cite{Israel1976}; the identification of modular flow with the boost for wedge algebras is Bisognano-Wichmann \cite{BisognanoWichmann1976}, with Sewell's black-hole reading \cite{Sewell1982}. Quasi-free states of the CAR algebra and their two-point calculus are Araki's \cite{Araki1971}; KMS theory, the ideal Fermi gas, and the inductive-limit results are standard in Bratteli-Robinson \cite{BratteliRobinson1997}. The geometric-modular-action programme, the candidate axiom for state selection, is Buchholz-Dreyer-Florig-Summers \cite{BDFS2000}. The Penrose transform is \cite{Penrose1969, EPW1981}; greybody filtering through the Regge-Wheeler and Teukolsky barriers is \cite{ReggeWheeler1957, Teukolsky1973, Page1976}. The boson-fermion correspondence and the charge-zero Fock isomorphism are classical two-dimensional field theory \cite{FrenkelKac1980}. Weinberg-Witten is \cite{WeinbergWitten1980} and the soft-graviton universality is \cite{Weinberg1964}. What this paper claims as new is the transport assembly: the I-state theorem at the operator level, all modes at once, with pairing and symbol both forced by the KMS identity (the single-mode occupancy is derived the same way in \cite{ButlerT5}); the exact finite-window continuum statements and the limit with its imports declared; the promotion of the homogeneity computation inside the proof of \cite{ButlerT9}'s dictionary theorem to a stated result the transform can consume; the identification of the lift's remaining joints with theorems of this series (the exterior exactness of \cite{ButlerT6}); and the quantified edge-support of the finite-mode corrections.

\subsection*{Contribution}

At the standard of \cite{ButlerT1}, with imports as Facts and the assembly joints itemised where a result is assembled.

\begin{enumerate}
\item The single-mode modular structure reproved self-contained, and the I-state theorem: existence of the quasi-free state whose modular clock is the geometry's clock, with gauge invariance and the Fermi-Dirac symbol both forced by the KMS identity, extending T5's single-mode derivation \cite[Thm.~7.20]{ButlerT5} to the full algebra, all modes at once; uniqueness beyond the quasi-free class as a corollary of the imported ideal-Fermi-gas theorem; the selection problem posed against T5's own analysis (Section \ref{sec:state}).
\item The continuum: exact finite-window KMS theorems, the inductive-limit assembly with its imports stated as Facts, free-fermion thermal field theory (Section \ref{sec:continuum}, Remark \ref{rem:tft}).
\item The lift: the Penrose homogeneity $-2s-2$ stated and proved from the twistor paper's ordering shift (Proposition \ref{prop:homogeneity}; the computation appears inside the proof of \cite[Thm.~5.7]{ButlerT9}, and is promoted here to a result the transform can consume), discharging the homogeneity half of the successor problem the twistor paper assigns to this paper, the transform itself staying open; the helicity lift and greybody inheritance assembled with every joint flagged, the interface premise (I) named in the statement (Section \ref{sec:lift}).
\item The $\ell_P^2/A$ correction: the exact edge identity for the finite window's current commutator, quantified tails, the area-suppressed reading, and the Weinberg-Witten triage (Section \ref{sec:seam}).
\item Machine verification of the computational identities in exact arithmetic \cite{ButlerVerifyT7}.
\end{enumerate}

\section{Setting and imports}\label{sec:setting}

Units are $G = c = \hbar = k_B = 1$ throughout, so for Schwarzschild the surface gravity is $\kappa = 1/4M$, the Hawking temperature is $\TH = \kappa/2\pi = 1/8\pi M$, the modular period is $\beta := 1/\TH = 8\pi M$, and areas are in Planck units, $\ell_P = 1$. (Two symbols are nonetheless kept visible where they carry the point: $\hbar$ in Section \ref{sec:lift}, where the ordering shift is a multiple of it, and $\ell_P$ in Section \ref{sec:seam}, where the corrections scale as $\ell_P^2/A$, $A$ the horizon area.)

\subsection{The single mode, self-contained}\label{subsec:mode}

The unit of the boundary mode space is a completed null pair: vectors $v, v^*$ with $v^2 = (v^*)^2 = 0$ and $vv^* + v^*v = 1$, the ladder pair the transport principle adjoins \cite{ButlerT1}. Everything modular about a single mode is elementary and is proved here outright, per the series' standard.

\begin{proposition}[Modular structure of the completed pair]\label{prop:mode}
Let $(v, v^*)$ be a completed pair and $\sigma := v^*v - vv^*$. Then:
\begin{enumerate}
\item $\sigma^2 = 1$; in the hyperbolic atom $\Cl(1,1)$ with $v = \tfrac12(e_0 + e_1)$, $v^* = \tfrac12(e_0 - e_1)$, one has $\sigma = e_0 e_1$: the unit boost bivector of the plane (relabelling the pair reverses the orientation);
\item the rotor flow $R_\lambda := \exp(\lambda\sigma/2)$ acts by $R_\lambda\, v\, R_\lambda^{-1} = e^{-\lambda}v$ and $R_\lambda\, v^*\, R_\lambda^{-1} = e^{+\lambda}v^*$: the pair diagonalises the boost with opposite weights;
\item the spectral projectors $\tfrac12(1 \pm \sigma)$ are the occupation projectors $v^*v$ and $vv^*$;
\item in the Fock representation the boost generator is the shifted number operator, $\tfrac12\sigma = N - \tfrac12$;
\item the bivector space of $\Cl(1,1)$ is one-dimensional, so the boost flow is the unique one-parameter rotor flow of the atom \emph{up to reparametrisation of the flow parameter and orientation}: the algebra fixes the orbits of any rotor-implemented modular flow, but not its speed, which is exactly the normalisation Remark \ref{rem:frontier} (Section \ref{sec:state}) leaves open.
\end{enumerate}
\end{proposition}

\begin{proof}
(1) $\sigma^2 = v^*vv^*v + vv^*vv^* - v^*v\,vv^* - vv^*\,v^*v = v^*v + vv^* - 0 - 0 = 1$, using $vv^*v = v$, $v^*vv^* = v^*$ and the vanishing squares. In $\Cl(1,1)$: $v^*v = \tfrac14(e_0 - e_1)(e_0 + e_1) = \tfrac14(e_0^2 - e_1^2 + e_0e_1 - e_1e_0) = \tfrac14(1 + 1 + 2e_0e_1) = \tfrac12(1 + e_0e_1)$, so $\sigma = v^*v - vv^* = e_0e_1$ in this orientation; relabelling the pair reverses it. (2) From $vv^*v = v$ and the vanishing squares, $\sigma v = (v^*v - vv^*)v = -vv^*v = -v$ and $v\sigma = v(v^*v - vv^*) = vv^*v = +v$. The first relation gives $e^{\lambda\sigma/2}\,v = v\,e^{-\lambda\sigma/2}$ term by term in the exponential series, so $R_\lambda\, v\, R_\lambda^{-1} = v\,e^{-\lambda\sigma}$. Since $\sigma^2 = 1$, $e^{-\lambda\sigma} = \cosh\lambda - \sigma\sinh\lambda$, and the second relation collapses it on the left factor: $v\,e^{-\lambda\sigma} = v(\cosh\lambda - \sinh\lambda) = e^{-\lambda}v$. The $v^*$ computation is the mirror image, with both signs reversed. (3) $\tfrac12(1 + \sigma) = \tfrac12(1 + v^*v - vv^*) = \tfrac12(v^*v + vv^* + v^*v - vv^*) = v^*v$, and likewise for the complement. (4) In the two-dimensional Fock module ($v^* = $ creation, $v = $ annihilation, $N = v^*v$): $\tfrac12\sigma = \tfrac12(v^*v - vv^*) = \tfrac12(N - (1 - N)) = N - \tfrac12$. (5) The bivector space of $\Cl(1,1)$ is the line $\R\,e_0e_1$, so every one-parameter rotor group is of the form $\lambda \mapsto \exp(c\lambda\sigma/2)$ for some real $c$: the same orbits at a rescaled parameter, which is the reparametrisation the statement reserves. The verification script checks every clause in exact arithmetic, including the finite conjugation of clause (2) at a rational rapidity.
\end{proof}

\begin{fact}[Imported: modular descent, its premises, and the geometric normalisation]\label{fact:descent}
Let the ascended algebra of a horizon pair carry a state that is pure, invariant under the boost flow of Proposition \ref{prop:mode}, and entangling across the pair (cyclic and separating for the one-sided subalgebra). Then the descended state is KMS for the boost flow at \emph{some} inverse temperature $\beta \in \R$ of rapidity, and the hypotheses leave $\beta$ free, its sign included; the Gibbs \emph{form} of the descended state is state-independent. Positivity and the value are supplied by T5's named premises: under (P) univalence and (R) the KMS condition at the geometric period (one modular period equal to one rapidity period $2\pi$, i.e.\ Schwarzschild time $2\pi/\kappa$, the Euclidean period $8\pi M$), the descended state is Fermi-Dirac at $\beta = 2\pi$ of rapidity, which is the physical temperature $\TH = \kappa/2\pi = 1/8\pi M$; the boost plane is the one the shape operator scales by $-2$, and $\kappa = 1/4M$ is proved from the gauge field. Per mode of Schwarzschild frequency $\omega$ the reduced density matrix is the Fermi-Dirac matrix $\mathrm{diag}(1, e^{-\omega/\TH})/Z$. \emph{Provenance:} T5 \cite[Thms.~7.3 and 7.6, Prems.~7.8 and 7.19, Prop.~7.18, Thm.~7.20, Cor.~6.4]{ButlerT5}, with $\kappa = 1/4M$ from the gauge field in \cite[App.~A]{ButlerT5}; premise (R) is proved irreducible there, the temperature being no algebra-intrinsic quantity \cite[Thm.~7.16]{ButlerT5}. The wedge-modular mechanism is Bisognano-Wichmann-Sewell \cite{BisognanoWichmann1976, Sewell1982}, and the thermofield reading is Israel's \cite{Israel1976}.
\end{fact}

\begin{fact}[Imported: the algebra of the boundary modes is the CAR algebra]\label{fact:car}
For even transport charge, the completed complexified limit of the tower of iterated completions is $\Cl(p,q)_\C \otimes \CAR$, the ground algebra tensored with the CAR algebra of countably many fermionic modes, and the mode operators are the adjoined null pairs themselves ($a_j = v_j$, $a_j^\dagger = v_j^*$, with the Fock adjointness). Here the \emph{transport charge} is the mod-8 invariant the tower's completions advance and the \emph{lineage} is the tower of iterated completions itself, both as in \cite{ButlerT3}; the physical lineage has charge six, so the theorem applies. \emph{Provenance:} the tower paper \cite[Prop.~4.1, Thms.~5.3 and 5.6]{ButlerT3}.
\end{fact}

\begin{fact}[Imported: the s-wave bosonic sector]\label{fact:boson}
Two halves, at two levels. \emph{Unconditional:} on a window of the boundary mode space, the number-conserving bilinears (currents) satisfy the free-boson commutation relations exactly in the completed limit, with an exact finite-window operator identity whose deviation is a difference of edge operators; and the bilinear subalgebra is globally invariant under the modular boost. \emph{Conditional on T5's Premise (Q), joint quasi-freeness with the s-wave weights:} the restriction of the Fermi-Dirac KMS state to the bilinear subalgebra is KMS at the same temperature (an instance of Lemma \ref{lem:restrict} below), giving the Bose-Einstein occupancy at $\TH$ for the s-wave sector; the linear weight assignment is the s-wave reduction's convention, adopted there as part of the premise. \emph{Provenance:} T5 \cite[Prem.~8.1, Thm.~8.2]{ButlerT5}; the correspondence itself is classical \cite{FrenkelKac1980}. The exact edge identity is reproved, in the form this paper needs, as Proposition \ref{prop:edge} below.
\end{fact}

\begin{fact}[Imported: the exterior is exactly general relativity]\label{fact:exterior}
On the region where the gauge field is invertible, and at $\Lambda = 0$, the series' description coincides exactly with classical general relativity: the field equations are equivalent to the Einstein equations, and the vacuum equations are exactly the covariant constancy of the scale tractor. \emph{Provenance:} \cite[Fact 2.1, Thms.~4.5 and 4.7]{ButlerT6} (Thm.~4.7 is the complementary $\Lambda \neq 0$ case), resting on Lasenby-Doran-Gull's equivalence theorem.
\end{fact}

\begin{fact}[Imported: the twistor storey, and the ordering shift]\label{fact:rent}
The spinor module of the completed algebra $\Cl(2,4)$ is twistor space $\T$, with $\dim_\C \T = 4$ (a \emph{storey} is the series' word for one level of the completion tower, as in \cite{ButlerT9}), the palatial commutation relation $[Z^\mu, \bar Z_\nu] = \hbar\,\delta^\mu_\nu$ on the quadratic (palatial) quantisation of the storey, and, for $M \in \mathfrak{gl}(\T)$, the symmetrised quadratic $q(M) := \tfrac12\bigl(Z^\mu M_\mu^{\ \nu}\bar Z_\nu + \bar Z_\nu M_\mu^{\ \nu} Z^\mu\bigr)$. (Twistor-spinor indices are written $\mu, \nu$ throughout; $\alpha$ and $\beta$ are reserved for the modular dynamics and its inverse temperature.) Symmetric ordering shifts the distinguished $\mathfrak{u}(1)$ generator $q(\id)$ by $-\hbar/2$ per complex dimension of $\T$, totalling $-2\hbar$ (equivalently, $q(M) = \sum M_\mu^{\ \nu} Z^\mu \bar Z_\nu - \tfrac{\hbar}{2}\tr M$); on the Clifford storey, the boost generator of each adjoined plane is $N - \tfrac12$, a shift of minus one half per fermionic mode: the two storeys carry the same signed ordering shift. \emph{Provenance:} the twistor paper \cite[Thm.~3.1, Rem.~3.2, \S 5, Lem.~5.5, Prop.~5.6, Thm.~5.7]{ButlerT9}; the homogeneity consequence is computed inside the proof of its Theorem 5.7 and is promoted to a stated result as Proposition \ref{prop:homogeneity} below.
\end{fact}

\subsection{The boundary-mode algebra, and its dynamics}\label{subsec:ledger}

\begin{definition}[The boundary-mode algebra, windows, and the boost dynamics]\label{def:ledger}
Fix a countable mode set $J$ with Schwarzschild frequencies $\omega_j > 0$ ($j \in J$), the frequencies the geometry assigns to the boundary modes (Fact \ref{fact:descent}). The \emph{boundary-mode algebra} $\mathfrak{A}$ is the CAR algebra over $\ell^2(J)$ (Fact \ref{fact:car}), the C*-inductive limit of the \emph{window algebras} $\mathfrak{A}_F$, $F \subset J$ finite, each $\mathfrak{A}_F \cong \End(\Lambda^\bullet \C^F)$ of dimension $4^{|F|}$. The \emph{boost dynamics} is the one-parameter group $\tau_t$ acting on generators by $\tau_t(a_j) = e^{-i\omega_j t}\,a_j$, with $t$ Schwarzschild time: the quasi-free flow whose one-particle Hamiltonian is $h = \mathrm{diag}(\omega_j)$; it preserves every window. In rapidity $\lambda = \kappa t$ the same flow has boost weights $\omega_j/\kappa$ and modular temperature $2\pi$ (Fact \ref{fact:descent}); in Schwarzschild time the inverse temperature is $\beta = 1/\TH = 2\pi/\kappa$, and this normalisation is used throughout.
\end{definition}

\section{The state input: the I-state}\label{sec:state}

T5 proves single-mode thermality from its premises (P) and (R), deriving the occupancy from the gluing identity, and says exactly what is premised. This section lifts the statement to the full mode algebra: within the quasi-free class the state is unique and its two-point data are forced outright, all modes at once; given the dynamics, even the class restriction can be removed, by a theorem of the literature; what cannot yet be derived is the dynamics-selection axiom itself, and that open problem is posed at the end, against T5's own selection analysis.

\subsection{Two elementary lemmas}\label{subsec:lemmas}

We write $\alpha_t$ for a finite-dimensional dynamics; on a window $F$ the relevant instance is $\alpha_t = \tau_t|_{\mathfrak{A}_F}$, generated by $H_F$, the second quantisation of $h|_F$.

\begin{lemma}[Gibbs states are KMS, and conversely, in finite dimensions]\label{lem:gibbs}
Let $\mathfrak{B} = \End(\C^m)$, let $H$ be diagonal in the basis $(\xi_a)$ with eigenvalues $\lambda_a$, all distinct or not, let $\alpha_t = \mathrm{Ad}(e^{itH})$, and $\beta \in \R$. Then the Gibbs state $\omega_G(x) := \tr(e^{-\beta H}x)/\tr(e^{-\beta H})$ satisfies the KMS identity
\[
\omega_G(x\,y) \;=\; \omega_G\bigl(y\,\alpha_{i\beta}(x)\bigr), \qquad \alpha_{i\beta}(x) := e^{-\beta H}x\,e^{\beta H},
\]
for all $x, y$; and it is the \emph{unique} $\alpha$-KMS state at $\beta$ on $\mathfrak{B}$, degenerate spectrum included.
\end{lemma}

\begin{proof}
The identity is cyclicity of the trace:
\[
\omega_G\bigl(y\,\alpha_{i\beta}(x)\bigr) = \tr\bigl(e^{-\beta H}y\,e^{-\beta H}x\,e^{\beta H}\bigr)/Z = \tr\bigl(e^{-\beta H}xy\bigr)/Z = \omega_G(xy),
\]
the trailing $e^{\beta H}$ moved around the trace to cancel the inner factor. Conversely write $\omega(x) = \tr(\rho x)$ and evaluate the KMS identity on matrix units $x = E_{ab}$, $y = E_{cd}$: since $\alpha_{i\beta}(E_{ab}) = e^{-\beta(\lambda_a - \lambda_b)}E_{ab}$ and $\tr(\rho\,E_{ab}E_{cd}) = \delta_{bc}\,\rho_{da}$, the identity reads
\[
\delta_{bc}\,\rho_{da} \;=\; e^{-\beta(\lambda_a - \lambda_b)}\,\delta_{da}\,\rho_{bc}
\qquad \text{for all } a, b, c, d.
\]
Taking $c = b$ and $d = a$ gives $\rho_{aa} = e^{-\beta(\lambda_a - \lambda_b)}\rho_{bb}$: Gibbs ratios between all diagonal entries. Taking $c = b$ and $d \neq a$ gives $\rho_{da} = 0$ (the right side carries $\delta_{da} = 0$), for every pair $d \neq a$, degenerate eigenvalues included: $\rho$ is diagonal. So $\rho = e^{-\beta H}/Z$, and the KMS state is unique. The verification script performs the solve symbolically for the mode algebra $\Mat_2$ with $H = \omega N$ and confirms both clauses as identities of rational functions in the Boltzmann factor.
\end{proof}

\begin{lemma}[KMS restricts to invariant subalgebras]\label{lem:restrict}
If $\omega$ is $\alpha$-KMS at $\beta$ on $\mathfrak{B}$ and $\mathfrak{C} \subseteq \mathfrak{B}$ is a unital subalgebra with $\alpha_t(\mathfrak{C}) = \mathfrak{C}$, then $\omega|_{\mathfrak{C}}$ is KMS at $\beta$ for $\alpha|_{\mathfrak{C}}$.
\end{lemma}

\begin{proof}
For $x, y \in \mathfrak{C}$ the correlation function $t \mapsto \omega(y\,\alpha_t(x))$ is computed entirely from elements of $\mathfrak{C}$ (invariance keeps $\alpha_t(x) \in \mathfrak{C}$), and is the same function of $t$ whether $x, y$ are regarded in $\mathfrak{C}$ or in $\mathfrak{B}$. Its continuation to the closed strip is unique (a function analytic on the strip and continuous on its closure is determined by its boundary restriction to the real line, by uniqueness of analytic continuation), so the strip function and the KMS boundary identity transfer verbatim from $\mathfrak{B}$ to $\mathfrak{C}$.
\end{proof}

\subsection{Quasi-free states, and the derived symbol}\label{subsec:quasifree}

\begin{fact}[Imported: quasi-free states of the CAR algebra]\label{fact:quasifree}
For every operator $A$ on $\ell^2(J)$ with $0 \le A \le 1$ (the \emph{symbol}) there is a unique gauge-invariant state $\omega_A$ on $\mathfrak{A}$ with two-point function $\omega_A(a^\dagger(f)\,a(g)) = \ip{g}{Af}$ and all higher truncated functions zero (Wick determinants). More generally, a quasi-free state (all truncated functions beyond second order vanishing) is determined by its full two-point data, the symbol together with the pairing $P(f, g) := \omega(a^\dagger(f)\,a^\dagger(g))$, and gauge invariance is exactly $P = 0$. On a finite window with diagonal symbol this is the explicit product state $\bigotimes_j \rho_j$, $\rho_j = \mathrm{diag}(1 - n_j,\, n_j)$ in the occupation basis, which the script realises exactly. \emph{Provenance:} Araki \cite{Araki1971}; \cite[\S 5.2]{BratteliRobinson1997}.
\end{fact}

\begin{fact}[Imported: KMS passes to inductive limits]\label{fact:limitkms}
Let a C*-algebra be the closure of an increasing union of subalgebras, each globally invariant under a one-parameter dynamics, and let a state be KMS at $\beta$ on each subalgebra for the restricted dynamics. Then the state is KMS at $\beta$ on the whole algebra: the analytic-element requirement is checked on the dense union and passes to the closure. \emph{Provenance:} \cite[Prop.~5.3.25]{BratteliRobinson1997}.
\end{fact}

\begin{fact}[Imported: quasi-free KMS reduces to the two-point level]\label{fact:twopoint}
For the CAR algebra over a one-particle space with the quasi-free dynamics of a self-adjoint one-particle Hamiltonian, a gauge-invariant quasi-free state is KMS at $\beta$ if and only if its two-point function satisfies the two-point KMS identity: through the Wick-determinant structure, the analytic-continuation requirement for products of generators reduces to the two-point level. \emph{Provenance:} Araki \cite{Araki1971}; \cite[\S 5.3]{BratteliRobinson1997}.
\end{fact}

\begin{theorem}[The I-state]\label{thm:istate}
Consider the boundary-mode algebra with its boost dynamics (Definition \ref{def:ledger}) at inverse temperature $\beta = 1/\TH = 2\pi/\kappa$.
\begin{enumerate}
\item \emph{(Existence, exactly on windows.)} The quasi-free state $\omega_I$ with the Fermi-Dirac symbol
\[
A_I \;=\; \bigl(1 + e^{\beta h}\bigr)^{-1}, \qquad\text{i.e.}\qquad n_j = \frac{1}{e^{\omega_j/\TH} + 1},
\]
restricts on every window to the Gibbs state of $H_F = \sum_{j \in F}\omega_j N_j$ and is KMS there, exactly (Lemma \ref{lem:gibbs}); it is the unique quasi-free extension of these window states, and it is $\tau$-KMS at $\beta$ on the full algebra by Fact \ref{fact:limitkms}.
\item \emph{(Gauge invariance is forced.)} If a quasi-free state of $\mathfrak{A}$ is $\tau$-KMS at $\beta$, its pairing vanishes: the state is gauge-invariant automatically.
\item \emph{(The symbol is derived.)} Its symbol then satisfies the operator identity $A(1 + e^{-\beta h}) = e^{-\beta h}$, equivalently $A = e^{-\beta h}(1 - A)$, whence $A = A_I$: within the quasi-free class, the state whose modular clock is the geometry's clock is unique, and its two-point data are forced.
\end{enumerate}
\end{theorem}

\begin{proof}
(1) On a window $F$, the quasi-free state with diagonal symbol is the product state of Fact \ref{fact:quasifree} with $n_j = (1 + e^{\beta\omega_j})^{-1}$, which is exactly the Gibbs density matrix of $H_F$: $e^{-\beta H_F} = \bigotimes_j \mathrm{diag}(1, e^{-\beta\omega_j})$, normalised; the script builds the two sides independently and confirms the identity in exact arithmetic. Lemma \ref{lem:gibbs} gives the KMS identity on the window, exactly; the script verifies it on a three-mode window at rational Boltzmann factors, on a spanning set of pairs. The windows are nested with dense union and the dynamics preserves each, so Fact \ref{fact:limitkms} passes the KMS property to $\mathfrak{A}$.
(2) Apply the KMS identity to the pair $x = a^\dagger(f)$, $y = a^\dagger(g)$, with $f, g$ supported in windows (so that $\alpha_{i\beta}(a^\dagger(f)) = a^\dagger(e^{-\beta h}f)$ is again a generator):
\[
P(f, g) \;=\; \omega\bigl(a^\dagger(f)\,a^\dagger(g)\bigr) \;=\; \omega\bigl(a^\dagger(g)\,a^\dagger(e^{-\beta h}f)\bigr) \;=\; P(g,\, e^{-\beta h}f).
\]
Antisymmetry of creation operators gives $P(f, g) = -P(g, f)$, so $P(g, (1 + e^{-\beta h})f) = 0$ for all $f, g$ in the dense span of the windows; $1 + e^{-\beta h}$ is bounded with bounded inverse and preserves that span, so $P = 0$.
(3) Apply the KMS identity to the pair $x = a^\dagger(f)$, $y = a(g)$:
\[
\omega\bigl(a^\dagger(f)\,a(g)\bigr) \;=\; \omega\bigl(a(g)\,a^\dagger(e^{-\beta h}f)\bigr)
\;=\; \ip{g}{e^{-\beta h}f} - \omega\bigl(a^\dagger(e^{-\beta h}f)\,a(g)\bigr),
\]
using the CAR $\{a(g), a^\dagger(f')\} = \ip{g}{f'}$. In symbols: $\ip{g}{Af} = \ip{g}{e^{-\beta h}f} - \ip{g}{Ae^{-\beta h}f}$ for all $f, g$ supported in windows, i.e.\ $A(1 + e^{-\beta h}) = e^{-\beta h}$ on the dense span, i.e.\ $A = e^{-\beta h}(1+e^{-\beta h})^{-1} = (1 + e^{\beta h})^{-1} = A_I$ (all four operators bounded, $h \ge 0$ here). By (2) and Fact \ref{fact:quasifree}, the state is determined by its two-point data, which are forced, so $\omega = \omega_I$. The mode-diagonal scalar identity $n = e^{-\beta\omega}(1-n)$, equivalently the detailed balance $n/(1-n) = e^{-\beta\omega}$, is checked by the script as an exact identity of rational functions.
\end{proof}

\begin{fact}[Imported: ideal-Fermi-gas uniqueness]\label{fact:fermigas}
For the CAR algebra over a one-particle space with self-adjoint $h$ and the induced quasi-free dynamics, the $\beta$-KMS state is unique and equals the quasi-free state with symbol $(1 + e^{\beta h})^{-1}$. \emph{Provenance:} \cite[Ex.~5.3.2 and \S 5.3]{BratteliRobinson1997}; the quasi-free calculus is \cite{Araki1971}.
\end{fact}

\begin{corollary}[Full uniqueness, given the dynamics]\label{cor:full}
Among all states of $\mathfrak{A}$, $\omega_I$ is the unique $\tau$-KMS state at $\beta$.
\end{corollary}

\begin{proof}
Fact \ref{fact:fermigas}, whose hypotheses Definition \ref{def:ledger} satisfies; the quasi-free state it names has symbol $A_I$, which is $\omega_I$ by Fact \ref{fact:quasifree}.
\end{proof}

\begin{remark}[What is now proved, and the open problem]\label{rem:frontier}
Three levels, stated separately. \emph{Proved here:} within the quasi-free class the state is unique, with gauge invariance and the Fermi-Dirac symbol both forced by the KMS identity (Theorem \ref{thm:istate}); in this sense the I-state is the framework's equilibrium vacuum, and the quasi-free restriction is a convenience of proof, removed at the next level. \emph{Imported:} given the boost dynamics and $\beta$, full uniqueness among all states (Corollary \ref{cor:full}, from Fact \ref{fact:fermigas}). \emph{Open, and named:} deriving the dynamics itself, i.e.\ an axiom selecting the boost flow and its temperature from the geometry. T5 has already treated this question at the level of a single factor: its geometric modular condition (G) is formulated and rejected there as over-determining the boundary modes, its premise (R) is proved irreducible, its rate-rigidity theorem shows the temperature is not algebra-intrinsic, and the resulting state class is recorded \cite[Def.~7.12, Thms.~7.16 and 7.28, Cor.~7.32]{ButlerT5}. The question this paper poses is different and remains open: a \emph{net-level} selection in the style of geometric modular action \cite{BDFS2000}, demanding that the modular groups of a net of boundary-mode algebras act geometrically and deriving the boost identification for the whole net rather than assuming it. Within the single atom the algebra already forces the orbits of any rotor-implemented flow, though not its speed (Proposition \ref{prop:mode}(5)); the execution for the net is a successor problem (Section \ref{sec:scope}). \emph{Point of contact:} the I-state on a horizon pair is Israel's thermofield double \cite{Israel1976}, and its GNS representation is the Araki-Wyss doubling; the non-equilibrium state of the descent, the cosmological state problem, is outside this paper's scope (successor problem 5).
\end{remark}

\section{The mode continuum}\label{sec:continuum}

The boundary mode space of Definition \ref{def:ledger} has countably many modes because the tower adjoins them one rung at a time (Fact \ref{fact:car}). The physical spectrum is emitted into a continuum of frequencies and a tower of angular channels, and this section carries the state theorem across that idealisation with the levels declared: what is exact, what is an inductive-limit result, and what is a modelling identification are separated.

\subsection{The continuum algebra}\label{subsec:contalg}

\begin{definition}[Continuum mode space]\label{def:continuum}
Let $\mathfrak{h}$ be a separable complex Hilbert space (the one-particle space; for the Schwarzschild s-wave sector, $L^2(\R_+, d\omega)$; with angular channels, a direct sum of copies indexed by $(\ell, m)$ and species), and let $h$ be self-adjoint on $\mathfrak{h}$ (the one-particle Hamiltonian of Schwarzschild time; for the s-wave sector, multiplication by the frequency $\omega$; self-adjointness, not positivity, is what the theorems below use, and Section \ref{sec:seam}'s bilateral ladder is an instance with spectrum unbounded below). The \emph{continuum mode space} is the CAR algebra $\CAR(\mathfrak{h})$ with the quasi-free (Bogoliubov) dynamics $\tau_t(a(f)) = a(e^{ith}f)$. The countable mode space of Definition \ref{def:ledger} is the case $\mathfrak{h} = \ell^2(J)$; the CAR functor turns isometric inclusions of one-particle spaces into unital embeddings of algebras, so refining the tower's mode set into a dense subset of the continuum realises $\CAR(\mathfrak{h})$ as the completed union of tower-type subalgebras: the continuum is the tower's limit re-parametrised, not a new object.
\end{definition}

\begin{theorem}[The I-state on the continuum]\label{thm:continuum}
For the continuum mode space at $\beta = 2\pi/\kappa$, with $h$ self-adjoint:
\begin{enumerate}
\item the quasi-free state $\omega_I$ with symbol $A_I = (1 + e^{\beta h})^{-1}$ is $\tau$-KMS at $\beta$: the two-point KMS identity is the computation of clause (3) run in reverse on the core, and Fact \ref{fact:twopoint} lifts it from the two-point level to the state. When $h$ has pure point spectrum, the finite spectral windows (spans of finitely many eigenvectors) are literal instances of Section \ref{sec:state}'s exact window theorem, proved outright; this is the case of the tower's mode space, Definition \ref{def:ledger};
\item the pairing of any quasi-free $\tau$-KMS state vanishes, by the computation of Theorem \ref{thm:istate}(2) performed on the core of clause (3);
\item the symbol derivation holds verbatim on the core: for $f$ in a bounded spectral subspace of $h$ (their union over cut-offs is a dense, $\tau$-invariant core on which $e^{-\beta h}f$ is defined and $a^\dagger(f)$ is entire analytic for $\tau$), the KMS identity forces $\ip{g}{A(1 + e^{-\beta h})f} = \ip{g}{e^{-\beta h}f}$, hence $A = A_I$ on the core and, $A$ and $A_I$ both bounded, everywhere: within the quasi-free class the continuum I-state is unique with derived two-point data. ($e^{-\beta h}$ may be unbounded when $h$ is unbounded below; the identity is only ever used on the core, and $A_I$ is bounded for every self-adjoint $h$.)
\end{enumerate}
Fact \ref{fact:fermigas} is stated for general self-adjoint $h$, so Corollary \ref{cor:full} also holds verbatim for the continuum mode space: among all states, $\omega_I$ is the unique $\tau$-KMS state.
\end{theorem}

\begin{proof}
(3) is the computation of Theorem \ref{thm:istate}(3) with the stated core in place of the window span: the core is $\tau$-invariant, $e^{-\beta h}$ preserves it and is bounded on it, and density plus boundedness of $A$ and $A_I$ extend the bounded identity $A = A_I$ from the core to $\mathfrak{h}$; the unbounded identity is used only on the core, where $1 + e^{-\beta h}$ acts bijectively. (2) is Theorem \ref{thm:istate}(2) on the same core. (1): the state $\omega_{A_I}$ exists for the symbol $0 \le A_I \le 1$ (Fact \ref{fact:quasifree}); its two-point function satisfies the two-point KMS identity by (3) read backwards, and Fact \ref{fact:twopoint} upgrades the two-point identity to the KMS property of the state. What this paper certifies independently, and the script checks exactly, is the window case, where no imported calculus is needed.
\end{proof}

\begin{fact}[Imported: the thermal representation]\label{fact:arakiwyss}
The GNS representation of $\CAR(\mathfrak{h})$ in the state $\omega_I$ is the Araki-Wyss representation: a doubled Fock space carrying two commuting copies of the CAR, with the GNS vector entangling them pairwise and the modular operator of the GNS vector implementing exactly the flow $\tau$ at inverse temperature $\beta$. \emph{Provenance:} Araki-Wyss \cite{ArakiWyss1964}; textbook form \cite[\S 5.3]{BratteliRobinson1997}.
\end{fact}

\begin{remark}[Free-fermion thermal field theory, and Israel's double]\label{rem:tft}
Fact \ref{fact:arakiwyss} is the continuum instance of T5's single-mode picture: the doubling is the completed pair structure mode by mode, the GNS vector is the thermofield double, and the modular flow is the boost. ``Free-fermion thermal field theory as the algebra carrying the mode continuum'' is now a theorem-backed description: the algebra is $\CAR(\mathfrak{h})$ (Fact \ref{fact:car} plus Definition \ref{def:continuum}), the state is the I-state (Theorem \ref{thm:continuum}), and the representation it generates is the standard thermal one. \emph{Point of contact:} Israel's thermofield reading of the eternal black hole \cite{Israel1976} is exactly this representation with the second copy read as the far wedge; T5's eviction dynamics runs inside it.
\end{remark}

\begin{remark}[Which dynamics, and at what level]\label{rem:whichdynamics}
The one-particle Hamiltonian $h$ was supplied by the geometry: the Schwarzschild frequency of each boundary mode, with the temperature's coefficient fixed by the surface gravity (Fact \ref{fact:descent}); in rapidity the weights are $\omega/\kappa$ at modular temperature $2\pi$. Two notes on scope. First, the greybody filtering of Section \ref{sec:lift} does \emph{not} deform $h$ or the KMS structure: transmission through the exterior barrier is a property of the propagation from the horizon to infinity, not of the horizon algebra; conflating the two is what the interface premise (I) of Assembly \ref{asm:greybody} guards against, and it is named in that assembly's statement. Second, the identification of the mode set (which $(\omega, \ell, m)$ channels, with which degeneracies) is a modelling input from the standard near-horizon reduction; the theorems of this section are exact for whatever separable $(\mathfrak{h}, h)$ that input delivers; the algebra supplies the theorems and the physics supplies the input.
\end{remark}

\section{The four-dimensional lift}\label{sec:lift}

The results so far are stated for the boundary modes, in the effectively two-dimensional near-horizon theory whose unit is the hyperbolic atom. The physical spectrum is four-dimensional radiation sorted by helicity and angular momentum and filtered by the exterior geometry. This section assembles the lift. It is the paper's principal assembly section, its Tier B, and the helicity assembly's joints are itemised by name so that each carries its own provenance: one joint is a classical theorem, one is a theorem of this series completed in Proposition \ref{prop:homogeneity} below, and one is modelling, flagged as such; the greybody assembly itemises three more and adds one named premise, (I), its only non-theorem joint. What the lift does and does not claim is bounded in Remark \ref{rem:liftscope}.

\subsection{Helicity is the charge of the twistor storey}\label{subsec:helicity}

\begin{fact}[Imported: the Penrose transform]\label{fact:penrose}
Massless free fields of helicity $s$ on (suitably convex regions of complexified) Minkowski space are in natural isomorphism with the first cohomology of the corresponding twistor-space region with coefficients in $\mathcal{O}(-2s-2)$: the homogeneity $-2s-2$ classifies the helicity-$s$ field. \emph{Provenance:} Penrose \cite{Penrose1969}; the sheaf-cohomological form \cite{EPW1981}.
\end{fact}

\begin{proposition}[The homogeneity $-2s-2$ is the ordering shift]\label{prop:homogeneity}
On twistor functions, realise the palatial relation of Fact \ref{fact:rent} by $\bar Z_\nu = -\hbar\,\partial/\partial Z^\nu$ (the realisation with the palatial sign $[Z^\mu, \bar Z_\nu] = +\hbar\,\delta^\mu_\nu$), let $\Upsilon := Z^\mu \partial_\mu$ be the homogeneity (Euler) operator, its eigenvalue written $\mathsf{h}$ in sans serif to keep it clear of the one-particle Hamiltonian $h$, and let the helicity operator be one half the distinguished charge, $\mathsf{s} := \tfrac12\, q(\id)$, the standard twistor helicity operator \cite{Penrose1969, EPW1981} in the symmetric ordering that Fact \ref{fact:rent} forces storey-wide. Then, as an operator identity,
\[
q(\id) \;=\; -\hbar\,(\Upsilon + 2),
\]
so a twistor function homogeneous of degree $\mathsf{h}$ is an eigenfunction of $\mathsf{s}$ with eigenvalue $s\hbar$ if and only if
\[
\mathsf{h} \;=\; -2s - 2.
\]
The $-2$ is exactly the ordering shift: without the symmetric-ordering shift of $-\hbar/2$ per complex dimension the offset would vanish. Combined with Fact \ref{fact:penrose}, the classification parameter of the Penrose transform is the completion's own $\mathfrak{u}(1)$ charge, its offset supplied by the ordering. The computation is not new to the series: it appears, without being stated as a result, inside the proof of \cite[Thm.~5.7]{ButlerT9}, where T9's script also checks it; what this proposition adds is the statement at result level, in the transform's context, discharging the homogeneity half of the successor problem T9 assigns to this paper (the transform itself stays open).
\end{proposition}

\begin{proof}
With $\bar Z_\nu = -\hbar\,\partial_\nu$,
\[
q(\id) \;=\; \tfrac12\sum_\mu\bigl(Z^\mu \bar Z_\mu + \bar Z_\mu Z^\mu\bigr)
\;=\; \sum_\mu Z^\mu \bar Z_\mu + \tfrac12\sum_\mu [\bar Z_\mu, Z^\mu]
\;=\; -\hbar\,\Upsilon - \tfrac12\cdot 4\hbar \;=\; -\hbar\,(\Upsilon + 2),
\]
the commutator term being $\sum_\mu[\bar Z_\mu, Z^\mu] = -4\hbar$ by the palatial relation, one $-\hbar$ per complex dimension of $\T$; halved by the symmetrisation, this is Fact \ref{fact:rent}'s shift of $-\hbar/2$ per complex dimension, totalling $-2\hbar$, and it is the only place the ordering enters. On a homogeneity-$\mathsf{h}$ eigenfunction, $\Upsilon f = \mathsf{h} f$, so $\mathsf{s} f = -\tfrac{\hbar}{2}(\mathsf{h} + 2)\,f$, and $\mathsf{s} f = s\hbar f$ if and only if $\mathsf{h} = -2s - 2$. The script verifies the displayed operator identity on all monomials of degree at most three in the four twistor variables, and the resulting homogeneity table at $s = \pm 1, \pm 2$.
\end{proof}

\begin{assembly}[The helicity lift, with its joints itemised]\label{asm:helicity}
Helicity is not additional structure imported into the framework: it is the eigenvalue of the distinguished $\mathfrak{u}(1)$ charge of the completion's twistor storey, and the four-dimensional mode space of each massless species has a canonical place in it. The joints, each with its provenance:
\begin{enumerate}
\item (\emph{Joint J1: the transform, classical.}) By Fact \ref{fact:penrose}, the helicity-$s$ mode space is the cohomology at homogeneity $-2s-2$; photon kinematics are at homogeneities $-4$ and $0$ ($s = \pm1$), graviton kinematics at $-6$ and $+2$ ($s = \pm2$).
\item (\emph{Joint J2: the charge, proved in the series.}) The ordering shift is proved in the twistor paper (Fact \ref{fact:rent}), and Proposition \ref{prop:homogeneity} derives the homogeneity $-2s-2$ from it: the classification parameter of Joint J1 is an internal charge of the completion, so the helicity ladder is part of the ascended kinematics rather than a phenomenological label.
\item (\emph{Joint J3: the reduction, modelling.}) The current boson of the s-wave theory (Fact \ref{fact:boson}) is the s-wave reduction of each such four-dimensional field, and the lift from the boundary modes to the field is the standard partial-wave reconstruction, each helicity sector identified through the transform of Joint J1: the near-horizon reduction of field theory on the Schwarzschild exterior identifies each $(\omega, \ell, m, s)$ channel with a two-dimensional mode, which is where Definition \ref{def:continuum}'s modelling input enters, flagged as such in Remark \ref{rem:whichdynamics}.
\end{enumerate}
The composite statement is the conjunction of the three joints with the provenance displayed: J1 is classical, J2 is Fact \ref{fact:rent} together with Proposition \ref{prop:homogeneity}, and J3 is the one modelling joint. The homogeneity table is checked by the script.
\end{assembly}

\subsection{Greybody inheritance}\label{subsec:greybody}

\begin{assembly}[Greybody inheritance, under the interface premise]\label{asm:greybody}
\emph{Premise (I), the interface identification:} the horizon-sector state feeds the exterior modes in the standard way, i.e.\ the occupancy of each near-horizon channel is carried to the corresponding asymptotic channel by the exterior propagation. This is the same assumption every derivation of Hawking flux makes, from \cite{Hawking1975} onward, and it is this assembly's one non-theorem joint. Under (I), the emitted flux in each channel $(\omega, \ell, m, s)$ equals (with $\mathcal{N}$ the count of emitted quanta)
\[
\frac{d\mathcal{N}}{dt\,d\omega} \;=\; \frac{1}{2\pi}\,\Gamma_{s\ell m}(\omega)\; n_\pm(\omega), \qquad
n_\pm(\omega) = \frac{1}{e^{\omega/\TH} \pm 1}
\]
(the upper sign for the fermionic species, the lower for the bosonic sector), with $\Gamma_{s\ell m}$ the classical transmission coefficient through the Regge-Wheeler or Teukolsky barrier. The joints: (i) the barriers and their transmission coefficients are properties of the exterior metric, and by Fact \ref{fact:exterior} the framework's exterior description \emph{is} classical general relativity on the invertible region, exactly and not asymptotically, so $\Gamma_{s\ell m}$ is inherited verbatim, mode by mode \cite{ReggeWheeler1957, Teukolsky1973}; (ii) the occupancies at the horizon are Theorem \ref{thm:continuum} for the fermionic species and Fact \ref{fact:boson}'s (Q)-conditional half for the bosonic sector, both at exactly $\TH$ with the temperature's coefficient fixed geometrically (Fact \ref{fact:descent}); (iii) given (I), the displayed formula is the definition of filtered thermal flux: the framework's flux prediction equals Hawking's \cite{Hawking1975, Page1976}, greybody structure included.
\end{assembly}

\begin{remark}[What the lift does and does not claim]\label{rem:liftscope}
The section's claim is deliberately bounded: the \emph{spectrum} of the four-dimensional radiation is derived (occupancies theorem-grade at the horizon under the named premises, filtering inherited exactly on the invertible region, the helicity location derived from the ordering shift). What is not claimed is a derivation of four-dimensional field \emph{dynamics} from the boundary modes: that is the task of the action paper T10 \cite{ButlerT10} and, for the emergent transverse geometry that would make Joint J3 itself a theorem, of the entropy paper T8 \cite{ButlerT8}, the fuzzy-sphere question being shared between T8 and T10. The lift beyond assembly level is a named successor problem (Section \ref{sec:scope}).
\end{remark}

\section{The finite boundary mode space, and the $\ell_P^2/A$ correction}\label{sec:seam}

The bosonisation bridge of T5 rests, at one joint, on the charge-zero Fock isomorphism of the \emph{infinite} chiral fermion \cite{FrenkelKac1980}, while the physical boundary mode space is finite: $L$ modes, with $L$ of order the horizon area $A$ in Planck units. This section proves that everything the finite window changes is supported at its edges, with the exact operator identity and quantified tails, and then records the reading: the corrections are area-suppressed, so the joint's finite-mode correction takes the form of a prediction.

\subsection{The exact edge identity}\label{subsec:edge}

Fix a window of $L$ fermionic modes $\psi_1, \dots, \psi_L$ (CAR: $\{\psi_k, \psi_l^\dagger\} = \delta_{kl}$, all other anticommutators zero) and the currents
\[
a_m \;:=\; \sum_{k=1}^{L-m} \psi_k^\dagger\,\psi_{k+m}, \qquad 1 \le m < L .
\]

\begin{proposition}[The current commutator, exactly, with its edge split]\label{prop:edge}
In the finite window, exactly and as an operator identity,
\[
[\,a_m,\, a_m^\dagger\,] \;=\; B_m - T_m, \qquad
B_m := \sum_{k=1}^{m} N_k, \qquad T_m := \sum_{k=L-m+1}^{L} N_k ,
\]
with $N_k = \psi_k^\dagger\psi_k$: the deviation from the infinite-volume central value is a difference of \emph{edge operators}, supported on the lowest and highest $m$ modes. Consequently, for $2m \le L$, on every state in which the bottom $m$ levels are filled and the top $m$ levels are empty (an \emph{edge-clean} state), $[a_m, a_m^\dagger]$ acts as the scalar $m$: the free-boson relation holds exactly, not asymptotically, on the edge-clean subspace. Moreover $[a_m, a_{m'}] = 0$ for all $m, m'$, exactly; and for $m \neq m'$ the cross commutator $[a_m, a_{m'}^\dagger]$ is a difference of bilinears supported on the lowest and highest $\max(m, m')$ modes, which annihilates every state that is edge-clean at depth $\max(m, m')$.
\end{proposition}

\begin{proof}
The fermionic bilinear commutator is $[\psi_a^\dagger\psi_b,\ \psi_c^\dagger\psi_d] = \delta_{bc}\,\psi_a^\dagger\psi_d - \delta_{ad}\,\psi_c^\dagger\psi_b$, an instance of the CAR. With $a_m^\dagger = \sum_{l=1}^{L-m}\psi_{l+m}^\dagger\psi_l$:
\[
[a_m, a_m^\dagger] = \sum_{k,l=1}^{L-m}\bigl(\delta_{k+m,\,l+m}\,\psi_k^\dagger\psi_l - \delta_{k l}\,\psi_{l+m}^\dagger\psi_{k+m}\bigr)
= \sum_{k=1}^{L-m} N_k - \sum_{k=m+1}^{L} N_k ,
\]
after shifting the index of the second sum. The shared middle range cancels pairwise: for $2m \le L$ the terms $k = m{+}1, \dots, L{-}m$ appear in both sums and cancel, leaving $B_m - T_m$; for $2m > L$ the two sums instead share the empty middle and the overlap $k = L{-}m{+}1, \dots, m$ of $B_m$ and $T_m$ cancels on the right-hand side, so the identity $\sum_{1}^{L-m} N_k - \sum_{m+1}^{L} N_k = B_m - T_m$ holds in both regimes. On an edge-clean state (which requires $2m \le L$ to be non-vacuous), $B_m$ acts as $m$ (each bottom mode occupied) and $T_m$ as $0$. For the cross commutators, the same computation gives
\[
[a_m, a_{m'}] = \sum_{k}\psi_k^\dagger\psi_{k+m+m'} - \sum_{k}\psi_k^\dagger\psi_{k+m+m'} = 0
\]
exactly, both surviving sums running over the same range $1 \le k \le L - m - m'$; and, for $m > m'$,
\[
[a_m, a_{m'}^\dagger] = \sum_{k=1}^{m'}\psi_k^\dagger\psi_{k+m-m'} \;-\; \sum_{k=L-m+1}^{L-m+m'}\psi_k^\dagger\psi_{k+m-m'},
\]
every surviving index within $\max(m,m')$ of an edge. Each bottom term annihilates an edge-clean state because its target level $k$ is occupied; each top term because its source level $k + m - m'$ lies in the empty top band. The case $m < m'$ follows by taking adjoints, $[a_m, a_{m'}^\dagger] = [a_{m'}, a_m^\dagger]^\dagger$, with the same band support; it annihilates edge-clean states by the mirrored argument (creation onto an occupied bottom level, lowering of an empty top level). The script verifies the identity on the full Fock space of an eight-mode window, state by state.
\end{proof}

\subsection{The Planck factor, and the window tails}\label{subsec:planck}

This subsection works on the bilateral ladder, $k \in \Z$ with equal level spacing $\epsilon$: the idealisation of a deep sea, an instance of Definition \ref{def:continuum} with $h = \epsilon\,\mathrm{diag}(k)_{k \in \Z}$, unbounded below. The currents of Section \ref{subsec:edge} extend to the bilateral sums $a_m = \sum_{k \in \Z}\psi_k^\dagger\psi_{k+m}$, read term by term inside I-state expectations, where the sums converge absolutely; the finite window $[-K, K]$ returns in clause (3).

\begin{proposition}[Bose-Einstein occupancy, with quantified window corrections]\label{prop:planck}
Let the ladder carry the linear s-wave weight assignment (the convention T5's Premise (Q) adopts, a modelling input here as there; Fact \ref{fact:boson}): level spacing $\epsilon$, sea at level $0$, I-state occupancies $n_k = (x^{-k} + 1)^{-1}$ with $x = e^{-\beta\epsilon} \in (0,1)$ the Boltzmann factor of the level spacing. Then:
\begin{enumerate}
\item \emph{(Detailed-balance identity.)} For every $k$ and $m \ge 1$, exactly,
\[
n_k - n_{k+m} \;=\; \bigl(x^{-m} - 1\bigr)\; n_{k+m}\,\bigl(1 - n_k\bigr).
\]
\item \emph{(The Planck factor, exactly, in the bilateral limit.)} Summing (1) over $k \in \Z$,
\[
\sum_{k \in \Z} n_{k+m}\,(1 - n_k) \;=\; \frac{m}{x^{-m} - 1} \;=\; \frac{m}{e^{\beta\epsilon m} - 1}:
\]
the current two-point function $\langle a_m^\dagger a_m\rangle_I = \sum_k n_{k+m}(1 - n_k)$ (a Wick computation) is exactly $m$ times the Bose-Einstein value at $\TH$, with no vacuum correction in this normalisation. The factor $m$ is the free-boson normalisation: on the bilateral ladder the currents are used only inside expectations, and $\langle [a_m, a_m^\dagger]\rangle_I = \sum_k (n_k - n_{k+m}) = m$, the telescoped left side of (1), matching the finite-window scalar of Proposition \ref{prop:edge}; the unit-normalised boson is $\hat a_m = a_m/\sqrt{m}$, whose occupancy $\langle \hat a_m^\dagger \hat a_m\rangle_I = 1/(e^{\beta\epsilon m} - 1)$ is the $n_-$ of the lift.
\item \emph{(Window tails.)} For the window $[-K, K]$ with the sea at the centre, i.e.\ the sum of (2) restricted to $k \in [-K, K-m]$ (both mode indices inside the window), the bilateral sum exceeds the window sum by a positive amount bounded by $2m\,x^{K - m}/(1 - x)$: exponentially small in the distance from the sea to the edges, in units of the level spacing.
\end{enumerate}
\end{proposition}

\begin{proof}
(1) With $u := x^{-k}$ and $w := x^{-m}$: $n_k - n_{k+m} = \tfrac{1}{u+1} - \tfrac{1}{uw+1} = \tfrac{u(w-1)}{(u+1)(uw+1)}$, and $(w-1)\,n_{k+m}(1-n_k) = (w-1)\cdot\tfrac{1}{uw+1}\cdot\tfrac{u}{u+1}$: identical. The script checks it as an identity of rational functions in the formal variable $u$ at several rational $w$.
(2) The terms $n_{k+m}(1 - n_k)$ decay like $x^{|k|}$ in both directions, so all sums converge absolutely and may be regrouped. Fix a residue class $k = r + jm$, $j \in \Z$: the class sum of the left side of (1) is $\lim_{j_0 \to \infty}\sum_{j=-j_0}^{j_0}\bigl(n_{r+jm} - n_{r+(j+1)m}\bigr) = \lim_{j_0 \to \infty}\bigl(n_{r-j_0 m} - n_{r+(j_0+1)m}\bigr) = 1 - 0$, the interior terms cancelling in adjacent pairs; there are $m$ residue classes, and dividing by $(x^{-m}-1)$ via (1) gives the display. The Wick computation: in the gauge-invariant quasi-free state, $\langle \psi^\dagger_{l+m}\psi_l\,\psi^\dagger_k\psi_{k+m}\rangle$ has vanishing disconnected part for $m \ge 1$ (the two-point function is mode-diagonal) and connected part $\delta_{kl}\,n_{k+m}(1 - n_k)$.
(3) Bound each term two ways, both valid for every integer index: $n_{k+m}(1-n_k) \le n_{k+m} \le x^{k+m}$ and $n_{k+m}(1-n_k) \le 1 - n_k \le x^{-k}$, since $n_j \le x^j$ and $1 - n_j \le x^{-j}$ are each the statement $1 \le 1 + x^{\pm j}$. The terms missing from the window sum are those with $k > K - m$ or $k < -K$, all positive, so the difference is positive and at most
\[
\sum_{k > K-m} x^{k+m} + \sum_{k < -K} x^{-k} \;=\; \frac{2\,x^{K+1}}{1-x} \;\le\; \frac{2m\,x^{K-m}}{1-x}.
\]
The script verifies instances of the bound exactly, and checks that the window sums at $K = 30$ and $K = 40$ agree with the closed form to within the proven tails, all in exact rational arithmetic.
\end{proof}

\begin{remark}[The $\ell_P^2/A$ correction, interpreted]\label{rem:seam}
Propositions \ref{prop:edge} and \ref{prop:planck} together are the quantitative form of T5's note on the level of its assembly joint. The charge-zero Fock isomorphism is exact for the infinite chiral fermion; for the physical boundary mode space of $L \sim A/\ell_P^2$ modes, every deviation is edge-supported (Proposition \ref{prop:edge}), and thermal expectations of the deviations carry explicit tails (Proposition \ref{prop:planck}(3)). The corrections therefore vanish as the area grows, at a rate set by the window geometry: the relative scale of generic finite-size effects is read dimensionally as $1/L = \ell_P^2/A$ (for the discretisation of mode sums and the Sugawara normalisation; proving that scaling per named observable is successor problem 4), while edge-clean bulk correlators at fixed $m$ are exponentially smaller, at the proven rates. The quantum-gravity corrections are area-suppressed and vanish where general relativity is well tested; any new physics in this derivation would appear in the $\ell_P^2/A$ correction, and that correction is a prediction of the derivation rather than a defect in it.
\end{remark}

\begin{remark}[Point of contact: the Weinberg-Witten triage]\label{rem:ww}
The bosonic sector's lift contains a graviton candidate (homogeneities $-6, +2$ in Assembly \ref{asm:helicity}), and a composite graviton must answer \cite{WeinbergWitten1980}: a theory with a Lorentz-covariant conserved stress tensor admits no massless spin-2 particle coupling to it. The evasion, recorded here as this paper's own claim and at Tier C: the theorem's hypotheses require the stress tensor and the spin-2 state to be defined on one and the same spacetime, and on the composite reading they are not; the covariant stress tensor of the boundary modes is defined in the two-dimensional reduced theory at the degenerate locus, the graviton candidate is a collective mode of the emergent four-dimensional description, and the hypotheses fail exactly where they fail for boundary gravitons in holography. The evidential triage, stated because the series states its own claims at their level: by \cite{Weinberg1964}, \emph{any} massless spin-2, composite or fundamental, couples universally to stress-energy, so the equivalence principle and the universal vertex are common ground and do not discriminate; what is distinctive to the composite reading is the ultraviolet behaviour and the $\ell_P^2/A$ correction of Remark \ref{rem:seam}, and the discriminating experiments, if any, are at the scale of that correction. Below the Planck scale no available experiment discriminates the composite reading from the fundamental one, so the reading is at present untestable; this paper says so rather than presenting the absence of a refutation as support. The engineering no-go and its ontology-independence are assigned to the action paper T10 \cite{ButlerT10}.
\end{remark}

\section{Scope, and successor problems}\label{sec:scope}

\subsection*{What is proved, what is imported, what is modelling}

The abstract states each claim at one of three tiers: Tier A for what this paper proves, Tier B for what it assembles from results proved elsewhere or from named premises, and Tier C for an interpretation offered as such. The first paragraph below is the Tier A content, Proposition \ref{prop:homogeneity} included, which the abstract meets as the lift's one proved joint; the second lists the imports that content rests on, certified by their source papers or the literature; the third is the assembly and modelling content, the assembly at Tier B and the modelling items stated as such, with the composite-graviton triage at Tier C. In the claim table of T17 \cite{ButlerT17} the paper as a whole is Tier B, because its physical conclusions pass through the assembly section, while its mathematical core is Tier A.

\emph{Proved in this paper, at the standard of \cite{ButlerT1}:} the single-mode modular structure (Proposition \ref{prop:mode}); the Gibbs/KMS equivalence with uniqueness in finite dimensions and the restriction lemma (Lemmas \ref{lem:gibbs}, \ref{lem:restrict}). Also proved here: the I-state theorem's window-exact existence (Theorem \ref{thm:istate}(1), window case), the forced gauge invariance, and the derivation of the Fermi-Dirac symbol from the KMS identity, on the boundary modes and on the continuum core (Theorems \ref{thm:istate}(2)--(3), \ref{thm:continuum}(2)--(3)); the homogeneity $-2s-2$ from the ordering shift (Proposition \ref{prop:homogeneity}); the exact edge identity, the edge-clean exactness of the free-boson relation, and the cross-commutator support bound (Proposition \ref{prop:edge}); the detailed-balance identity, the exact bilateral Planck factor, and the window tail bounds (Proposition \ref{prop:planck}, on the (Q)-convention ladder named in its statement).

\emph{Imported as Facts with provenance:} T5's modular descent under its premises (P) and (R), with the geometric normalisation (Fact \ref{fact:descent}); the tower's CAR limit (Fact \ref{fact:car}); the s-wave bosonisation results, their occupancy half under T5's premise (Q) (Fact \ref{fact:boson}); the exterior exactness at $\Lambda = 0$ (Fact \ref{fact:exterior}); the twistor storey and the ordering shift (Fact \ref{fact:rent}). From the literature: Araki's quasi-free calculus, the inductive-limit KMS result, the two-point reduction, and the ideal-Fermi-gas uniqueness (Facts \ref{fact:quasifree}, \ref{fact:limitkms}, \ref{fact:twopoint}, \ref{fact:fermigas}); the Araki-Wyss representation (Fact \ref{fact:arakiwyss}); the Penrose transform (Fact \ref{fact:penrose}); the soft-graviton universality \cite{Weinberg1964} used in Remark \ref{rem:ww}.

\emph{Assembly and modelling, flagged at their level:} the helicity lift's joint J3 and the greybody interface premise (I) (Section \ref{sec:lift}); the identification of the physical mode set, and the linear s-wave weight convention of the bilateral ladder (Remark \ref{rem:whichdynamics}, Proposition \ref{prop:planck}); the $\ell_P^2/A$ dimensional reading of generic finite-size corrections (Remark \ref{rem:seam}); the composite-graviton reading and its triage (Remark \ref{rem:ww}), which is Tier C and says so.

\subsection*{Successor problems}

Six problems for the series' list of open directions.

\begin{enumerate}
\item \emph{State selection from geometry.} Formulate the geometric-modular-action axiom for a net of boundary-mode algebras over the horizon and derive the boost dynamics, the temperature entering through the net's geometric normalisation rather than as an algebra-intrinsic quantity, which T5's rate-rigidity theorem shows it cannot be \cite[Thm.~7.16]{ButlerT5}. Proposition \ref{prop:mode}(5) fixes the single-atom orbits, and T5's factor-level analysis (its condition (G), rejected as over-determining) records what is already established.
\item \emph{The lift beyond assembly.} Derive joint J3, the partial-wave identification, from the emergent transverse geometry (the fuzzy sphere of the counting programme), converting the helicity lift (Assembly \ref{asm:helicity}) from assembly to theorem; shared with the entropy paper T8 \cite{ButlerT8} and the action paper T10 \cite{ButlerT10}.
\item \emph{The Kerr continuum.} Superradiant channels make the rotating analogue of Theorem \ref{thm:continuum} genuinely different (no global KMS state at fixed angular velocity for all modes); formulate the correct rotating statement.
\item \emph{The $\ell_P^2/A$ correction, per observable.} Refine Remark \ref{rem:seam} into a table of leading finite-$L$ corrections for named observables (spectrum distortion, correlator deficits), with the $\ell_P^2/A$ scaling proved rather than dimensional.
\item \emph{The cosmological state.} Connect the equilibrium I-state to the non-equilibrium state of the descent.
\item \emph{Helicity counting per boundary mode.} Decide whether one boundary mode of the near-horizon theory carries one helicity channel of the lift or both helicities of a massless species. The two answers differ downstream by a factor of two in the capacity per boundary mode ($\ln 2$ against $2\ln 2$) and so in the entropy paper's calibration residue, $\kappa_c = 1/(4\ln 2)$ against $1/(8\ln 2)$ \cite[Prop.~4.12]{ButlerT8}. In this paper the question sits at joint J3 of Assembly \ref{asm:helicity}, the partial-wave identification of each $(\omega, \ell, m, s)$ channel with a two-dimensional mode, which is modelling; joints J1 and J2 place the two helicities at distinct homogeneities but do not say how many channels one boundary mode registers. T8 takes the one-helicity reading, at the tier of that joint; settling it is part of converting J3 to a theorem (problem 2).
\end{enumerate}

\section{Position in the series}\label{sec:position}

This paper is T7 of the series. \emph{It relies on six predecessors:} the primer T0 \cite{ButlerT0} and the root paper T1 \cite{ButlerT1} for the completed-pair language and the transport principle; the tower paper T3 \cite{ButlerT3} for the CAR limit that makes the boundary modes an algebra (Fact \ref{fact:car}); the singularity paper T5 \cite{ButlerT5} for modular descent under its premises (P) and (R), the bosonisation bridge under (Q), and the surface gravity $\kappa = 1/4M$ (Facts \ref{fact:descent} and \ref{fact:boson}); the interior paper T6 \cite{ButlerT6} for the exterior exactness that greybody inheritance rests on (Fact \ref{fact:exterior}); and the twistor paper T9 \cite{ButlerT9} for the ordering shift from which Proposition \ref{prop:homogeneity} is promoted.

\emph{Papers that rely on this one:} the entropy paper T8 \cite{ButlerT8} imports the single-mode modular structure, the quasi-free calculus, and the finite-window identities; the action paper T10 \cite{ButlerT10} imports the finite-window quantification of the current algebra; the standard-model paper T12 \cite{ButlerT12} imports the helicity identification of Proposition \ref{prop:homogeneity}; and the capstone T17 \cite{ButlerT17} tabulates this paper's claims at their tiers.

\section{Machine verification}\label{sec:verification}

The computational identities asserted in this paper are verified by the script deposited with it, \texttt{verify\_T07.py} \cite{ButlerVerifyT7}. The script imports the series' shared exact-arithmetic library \texttt{GA.py}, runs entirely in exact arithmetic (rational numbers and rational functions in a formal Boltzmann variable, no floating point anywhere), and so certifies each check as an identity, or as an exactly evaluated inequality, rather than an approximation. Twenty-five checks pass. Not scripted, and stated as such where they occur: the imported Facts, certified by their source papers' scripts or by the literature; the continuum passage beyond its window content; the assembly joints of Section \ref{sec:lift}; the classical Penrose transform; and the geometric inputs $\kappa = 1/4M$ and $\TH = \kappa/2\pi$, proved and scripted in T5.

The coverage, section by section.

\begin{description}
\item[Section \ref{sec:setting}, the single mode (Proposition \ref{prop:mode}).] All five clauses in the basis-blade arithmetic over $\Cl(1,1)$: the null pair's squares and anticommutator, $\sigma = e_0e_1$ with $\sigma^2 = 1$, the one-sided weight relations, the finite conjugation at the rational rotor $R = \tfrac53 + \tfrac43\,\sigma$ (that is, $\cosh(\lambda/2) = \tfrac53$) with scale factors $\tfrac19$ and $9$, the projector identities, and the one-dimensionality of the bivector space; plus the Fock identity $\tfrac12\sigma = N - \tfrac12$.
\item[Section \ref{sec:state}, the state (Lemma \ref{lem:gibbs}, Theorem \ref{thm:istate}).] The Gibbs KMS identity on all sixteen matrix-unit pairs of the mode algebra, as identities of rational functions in the formal Boltzmann factor; the uniqueness solve, performed symbolically, the KMS constraints on an arbitrary density matrix having a one-dimensional kernel spanned by the Gibbs matrix. Then the symbol itself: the Fermi-Dirac detailed-balance identity $n/(1-n) = e^{-\beta\omega}$; the Fermi-Dirac product state on a three-mode window confirmed equal to the Gibbs matrix of $H_F$, the two sides built independently, then confirmed KMS on all $4096$ matrix-unit pairs; and, on a two-mode window, the full KMS constraint system solved for an arbitrary density matrix against independently supplied dynamics, with one-dimensional kernel, vanishing pairing, and the Fermi-Dirac symbol.
\item[Section \ref{sec:continuum}, the continuum.] No separately computable identity: its exact content is the window case, scripted under Section 3, and its remaining steps are the imported calculus of Facts \ref{fact:twopoint} and \ref{fact:fermigas}, certified by their sources.
\item[Section \ref{sec:lift}, the lift (Proposition \ref{prop:homogeneity}).] The operator identity $q(\id) = -\hbar(\Upsilon + 2)$ on all twistor monomials of degree at most three, and the homogeneity table $-2s-2 \mapsto (-4, 0, -6, +2)$ at $s = \pm 1, \pm 2$.
\item[Section \ref{sec:seam}, the finite window (Propositions \ref{prop:edge} and \ref{prop:planck}).] The edge identity $[a_m, a_m^\dagger] = B_m - T_m$ as an operator identity on all $256$ states of an eight-mode window for $m = 1, 2, 3$; the exact free-boson relation on every edge-clean state; the exact vanishing of $[a_m, a_{m'}]$ and the vanishing of the cross commutator on edge-clean states. Then the ladder: the detailed-balance identity of Proposition \ref{prop:planck}(1) in the formal $x^{-k}$ at three values of $x^{-m}$, with the Bose-Einstein form $\bar n/(1 + \bar n) = e^{-\beta\omega}$ of the clause-(2) occupancy, at $\omega = \epsilon m$ on the ladder; the window sums' approach to the exact Planck factor from below, within the proven tail bound, as exact rational inequalities at two temperatures and three $m$; and the Wick formula checked against the direct density-matrix expectation on the full $256$-state window, the ladder in the paper's occupancy convention.
\end{description}

The script exits nonzero on any failure and prints one line per check, tagged by the statement it certifies.

\begin{thebibliography}{99}

\bibitem{ButlerT1}
L.~Butler,
\emph{The Transport Principle: Total Division in Clifford Algebras via Minimal Completion},
Transport Papers T1 (2026).

\bibitem{ButlerT3}
L.~Butler,
\emph{The Tower and the CAR Limit: Iterated Transport, the Frozen Clock, and Second Quantisation},
Transport Papers T3 (2026).

\bibitem{ButlerT5}
L.~Butler,
\emph{Division by Zero at the Gravitational Singularity: The Transport Principle Applied to the Black Hole, and the Information Paradox as a Category Error},
Transport Papers T5 (2026).

\bibitem{ButlerT6}
L.~Butler,
\emph{The Geometry of the Black-Hole Interior: Where Division Fails, the Weighted Limit, and the Kerr Tests},
Transport Papers T6 (2026).

\bibitem{ButlerT9}
L.~Butler,
\emph{The Twistor Correspondence: Ascended Spinors, Incidence, and the Palatial Bridge},
Transport Papers T9 (2026).

\bibitem{ButlerT0}
L.~Butler,
\emph{An Introduction to Geometric Algebra: A Primer for the Transport Papers},
Transport Papers T0 (2026).

\bibitem{ButlerT8}
L.~Butler,
\emph{Black-Hole Entropy and Scrambling: The Singularity Ledger and the No-Echo Prediction},
Transport Papers T8 (2026).

\bibitem{ButlerT10}
L.~Butler,
\emph{Dark Energy and the Action of Gravity: MacDowell-Mansouri, the Frozen Point, and One Way to Grip},
Transport Papers T10 (2026).

\bibitem{ButlerT12}
L.~Butler,
\emph{The Standard Model from the Minimal Massive Composite},
Transport Papers T12 (2026).

\bibitem{ButlerT17}
L.~Butler,
\emph{The Transport Programme: Method, Status, and the Falsifiable Ledger},
Transport Papers T17 (2026).

\bibitem{Hawking1975}
S.~W.~Hawking,
Particle creation by black holes,
\emph{Commun.\ Math.\ Phys.} \textbf{43} (1975), 199--220.

\bibitem{Israel1976}
W.~Israel,
Thermo-field dynamics of black holes,
\emph{Phys.\ Lett.\ A} \textbf{57} (1976), 107--110.

\bibitem{BisognanoWichmann1976}
J.~J.~Bisognano and E.~H.~Wichmann,
On the duality condition for quantum fields,
\emph{J.\ Math.\ Phys.} \textbf{17} (1976), 303--321.

\bibitem{Sewell1982}
G.~L.~Sewell,
Quantum fields on manifolds: PCT and gravitationally induced thermal states,
\emph{Ann.\ Phys.} \textbf{141} (1982), 201--224.

\bibitem{Araki1971}
H.~Araki,
On quasifree states of CAR and Bogoliubov automorphisms,
\emph{Publ.\ RIMS Kyoto Univ.} \textbf{6} (1971), 385--442.

\bibitem{ArakiWyss1964}
H.~Araki and W.~Wyss,
Representations of canonical anticommutation relations,
\emph{Helv.\ Phys.\ Acta} \textbf{37} (1964), 136--159.

\bibitem{BratteliRobinson1997}
O.~Bratteli and D.~W.~Robinson,
\emph{Operator Algebras and Quantum Statistical Mechanics 2}, 2nd ed.,
Springer, Berlin (1997).

\bibitem{BDFS2000}
D.~Buchholz, O.~Dreyer, M.~Florig and S.~J.~Summers,
Geometric modular action and spacetime symmetry groups,
\emph{Rev.\ Math.\ Phys.} \textbf{12} (2000), 475--560.

\bibitem{Penrose1969}
R.~Penrose,
Solutions of the zero-rest-mass equations,
\emph{J.\ Math.\ Phys.} \textbf{10} (1969), 38--39.

\bibitem{EPW1981}
M.~G.~Eastwood, R.~Penrose and R.~O.~Wells,
Cohomology and massless fields,
\emph{Commun.\ Math.\ Phys.} \textbf{78} (1981), 305--351.

\bibitem{ReggeWheeler1957}
T.~Regge and J.~A.~Wheeler,
Stability of a Schwarzschild singularity,
\emph{Phys.\ Rev.} \textbf{108} (1957), 1063--1069.

\bibitem{Teukolsky1973}
S.~A.~Teukolsky,
Perturbations of a rotating black hole I,
\emph{Astrophys.\ J.} \textbf{185} (1973), 635--647.

\bibitem{Page1976}
D.~N.~Page,
Particle emission rates from a black hole,
\emph{Phys.\ Rev.\ D} \textbf{13} (1976), 198--206.

\bibitem{FrenkelKac1980}
I.~B.~Frenkel and V.~G.~Kac,
Basic representations of affine Lie algebras and dual resonance models,
\emph{Invent.\ Math.} \textbf{62} (1980), 23--66.

\bibitem{WeinbergWitten1980}
S.~Weinberg and E.~Witten,
Limits on massless particles,
\emph{Phys.\ Lett.\ B} \textbf{96} (1980), 59--62.

\bibitem{Weinberg1964}
S.~Weinberg,
Photons and gravitons in S-matrix theory,
\emph{Phys.\ Rev.} \textbf{135} (1964), B1049--B1056.

\bibitem{ButlerVerifyT7}
L.~Butler,
\emph{verify\_T07.py: exact-arithmetic verification of the computational identities of ``Hawking Radiation from the Ledger''},
Python script, deposited with this paper alongside the shared library GA.py (2026).

\end{thebibliography}

\end{document}
